\chapter{I design pattern}
% Il codice in linea (\code) non va a capo: un po' di elasticita' in piu' evita righe che sbordano
\emergencystretch=2.5em

Nei capitoli sulle architetture abbiamo deciso la struttura \textbf{ad alto livello} di un sistema: quali sottosistemi,
quali livelli, come comunicano. Ogni sottosistema, però, è fatto di molte classi, e anche a questo livello più basso gli
stessi problemi si ripresentano da un progetto all'altro: creare un oggetto complicato, far collaborare due classi nate
separatamente, avvisare automaticamente chi è interessato a un cambiamento. Questo capitolo raccoglie la lezione
\textit{Software Design: Design Patterns}: che cos'è un design pattern, da dove viene, e i pattern più importanti, ognuno con
il suo codice Java (compilato ed eseguito: gli output riportati sono quelli veri) e, dove aiuta, con un \textbf{disegno della
memoria} che mostra che cosa succede davvero agli oggetti.

% ============================================================
\section{Dove siamo nel processo}

\waterfallbar{3}

I design pattern appartengono alla terza fase della cascata, \textit{Software e UI Design}, presentata nel Capitolo~3. Il
System Design ha definito l'architettura e assegnato i requisiti ai sottosistemi; ora, per ciascun sottosistema, si
definiscono gli \textbf{oggetti} necessari a realizzarlo. Gli strumenti della fase sono i \textbf{design pattern}
(questo capitolo), l'\textit{usability engineering} e i \textit{wireframe} ad alta fedeltà dell'interfaccia utente.

\subsection{Dalle architetture ai design pattern}

\begin{itemize}
  \item Un'architettura definisce la struttura di un'applicazione \textbf{ad alto livello}: in ogni modulo si raggruppano
        molte classi.
  \item Grazie alle architetture già note (client/server, a livelli, a microservizi\ldots) abbiamo una vasta scelta di
        soluzioni \textbf{ampiamente collaudate} per una grande classe di problemi.
  \item Esiste qualcosa di simile a un livello più basso? Esistono modi «standard» di combinare le classi fra loro per
        svolgere funzioni tipiche? Sì: i \textbf{design pattern}.
\end{itemize}

\Needspace{14\baselineskip}
\begin{center}
\begin{tikzpicture}
  % l'architettura: tre livelli
  \node[font=\scriptsize\bfseries, text=ph2!80!black] at (0,2.15) {Architettura (Capitoli 11 e 12)};
  \node[card=ph2, minimum width=3.4cm, minimum height=0.8cm] (l1) at (0,1.3) {Presentazione};
  \node[card=ph2, minimum width=3.4cm, minimum height=0.8cm] (l2) at (0,0) {Logica applicativa};
  \node[card=ph2, minimum width=3.4cm, minimum height=0.8cm] (l3) at (0,-1.3) {Accesso ai dati};
  \draw[flow] (l1) -- (l2);
  \draw[flow] (l2) -- (l3);
  % lo zoom su un sottosistema
  \draw[ph3!70, line width=0.7pt, dashed] (l2.north east) -- (3.6,1.9);
  \draw[ph3!70, line width=0.7pt, dashed] (l2.south east) -- (3.6,-1.9);
  \draw[ph3!70, line width=0.9pt, rounded corners=4pt, fill=ph3!4] (3.6,-1.9) rectangle (13.6,1.9);
  \node[font=\scriptsize\bfseries, text=ph3!80!black, anchor=north west] at (3.7,1.85)
    {Dentro un sottosistema: un design pattern (questo capitolo)};
  \umlclass[2.4cm]{dps}{(5.4,0.35)}{Soggetto}{}{attach(o)\\notify()}
  \umlclass[2.3cm]{dpo}{(9.7,0.35)}{\textit{Osservatore}}{}{update()}
  \umlclass[2.3cm]{dpc}{(9.7,-1.3)}{Tabella}{}{update()}
  \draw[navig] (dps.east) -- node[molt, above]{osservatori} node[molt, below, pos=0.82]{0..*} (dpo.west);
  \draw[realizza] (dpc.north) -- (dpo.south);
  \node[note, anchor=west, text width=2.3cm] at (11.15,0.35) {poche classi che collaborano in un modo collaudato};
\end{tikzpicture}
\end{center}

% ============================================================
\section{Che cos'è un design pattern}

\dfn{Design pattern}{
  Un \term{design pattern} è una \textbf{soluzione comune a un problema ricorrente} di progettazione. Astrae una struttura
  di progetto che si ripresenta in molti sistemi e la descrive in termini di classi e oggetti: le loro
  \textbf{dipendenze}, le loro \textbf{strutture}, le loro \textbf{interazioni} e le \textbf{convenzioni} da rispettare.
  Dà un \textbf{nome} alla struttura, la specifica in modo esplicito e distilla l'esperienza di chi l'ha già usata.
}

\subsection{Un po' di storia}

\begin{itemize}
  \item L'idea viene dall'architettura, quella degli edifici. L'architetto \textbf{Christopher Alexander} raccolse in
        \textit{A Pattern Language: Towns, Buildings, Construction} (1977) 253 pattern per progettare città e case, come
        «luce da due lati in ogni stanza»: le stanze illuminate da un lato solo sono sgradevoli, e la soluzione si può
        riusare in mille case diverse. Ogni pattern descrive un problema che si ripete e il nucleo della sua soluzione, da
        adattare ogni volta.
  \item Nel 1995 la «\textbf{Gang of Four}» (GoF: Erich Gamma, Richard Helm, Ralph Johnson e John Vlissides) portò l'idea
        nel software con il libro \textit{Design Patterns: Elements of Reusable Object-Oriented Software}, che descrive 23
        pattern per la programmazione a oggetti. Per Java si può vedere anche S.~J.~Metsker, \textit{Design Patterns in
        Java}, Pearson, 2003.
  \item Da allora ne sono stati proposti molti altri, in conferenze, simposi e libri. Non è raro che un singolo
        sviluppatore proponga un nuovo pattern e lo diffonda in modo informale, con blog e riviste.
\end{itemize}

\subsection{Le quattro parti di un pattern}

\begin{table}[H]
\centering
\renewcommand{\arraystretch}{1.35}
\begin{tabularx}{\textwidth}{>{\bfseries}P{3.2cm} L}
\toprule
Parte & Che cosa contiene \\
\midrule
1. Nome & Una o due parole che identificano il pattern (Singleton, Observer\ldots) e diventano parte del vocabolario dei
          progettisti. \\
2. Problema & Quando applicarlo: il problema da risolvere e il contesto in cui si presenta. \\
3. Soluzione & Gli elementi del progetto (classi e oggetti), le loro relazioni, responsabilità e collaborazioni. Non è un
               progetto concreto ma uno \textbf{schema}, da riempire ogni volta. \\
4. Conseguenze & I risultati e i compromessi (\textit{trade-off}) dell'applicazione: che cosa si guadagna e che cosa si
                 paga. \\
\bottomrule
\end{tabularx}
\caption{Le quattro parti di un design pattern.}
\end{table}

Un design pattern è inoltre:
\begin{itemize}
  \item \textbf{indipendente dal linguaggio e dall'implementazione}: descrive un'idea, che si può realizzare in Java, in
        C\# o in Python;
  \item una \textbf{micro-architettura}: poche classi che collaborano, non un sistema intero;
  \item da applicare \textbf{con giudizio}, mai meccanicamente: la soluzione va tradotta in termini concreti nel contesto
        dell'applicazione, ed è compito dello sviluppatore farlo.
\end{itemize}

\subsection{A che cosa servono}

\begin{itemize}
  \item \textbf{Codificare la buona progettazione}: distillano e diffondono l'esperienza, aiutano allo stesso modo
        principianti ed esperti, insegnano un modo di ragionare sul progetto.
  \item \textbf{Dare nomi espliciti alle strutture}: un vocabolario comune riduce la complessità e aumenta
        l'espressività. «Qui ci vuole un Observer» comunica in quattro parole una struttura di quattro classi.
  \item \textbf{Catturare e conservare le informazioni di progetto}: le decisioni si descrivono in modo sintetico e la
        documentazione migliora.
  \item \textbf{Facilitare la ristrutturazione} (\textit{refactoring}): i pattern sono collegati fra loro e danno al codice
        una flessibilità in più.
\end{itemize}

% ============================================================
\section{Come leggere i disegni della memoria}

Molti pattern si capiscono davvero solo guardando gli \textbf{oggetti} mentre il programma gira: quanti ne esistono, chi
punta a chi, che cosa resta in memoria. Per questo, nel capitolo, accanto al codice compaiono dei \textbf{disegni della
memoria} di un programma Java, divisi in tre zone.

\begin{itemize}
  \item \textbf{Stack}. Ogni chiamata di metodo in corso ha un suo \textit{frame}, che contiene i parametri e le variabili
        locali. Il frame nasce quando il metodo viene chiamato e sparisce quando il metodo termina; il metodo in esecuzione
        è quello disegnato più in alto.
  \item \textbf{Heap}. Qui vivono gli \textbf{oggetti}, creati con \code{new}. Un oggetto resta in memoria finché è
        \textbf{raggiungibile}, cioè finché una catena di riferimenti, partendo dallo stack o dai campi statici, arriva fino a
        lui; quando non lo è più, il \textit{garbage collector} lo elimina. Nei disegni gli oggetti irraggiungibili hanno il
        bordo tratteggiato.
  \item \textbf{Campi statici}. I campi \code{static} appartengono alla \textbf{classe}, non a un oggetto: ne esiste una sola
        copia, creata quando la classe viene caricata. (Nella JVM più diffusa, HotSpot, anche questi valori stanno
        tecnicamente nello heap, dentro l'oggetto che rappresenta la classe; li disegniamo a parte perché ne esiste una
        copia sola.)
\end{itemize}

Una variabile di tipo primitivo (\code{int}, \code{double}, \code{boolean}) contiene direttamente il suo valore. Una
variabile di tipo oggetto contiene invece un \textbf{riferimento}, disegnato come un pallino con una freccia verso
l'oggetto; \code{null} è il riferimento che non punta a nulla. Anche le stringhe sono oggetti, ma per non affollare i
disegni le scriviamo direttamente nel campo, fra virgolette.

\begin{esempio}{Due riferimenti, un solo oggetto}
\noindent\begin{minipage}[c]{0.6\textwidth}
\begin{lstlisting}[style=java, name=Punto.java]
public class Punto {
    int x, y;

    Punto(int x, int y) {
        this.x = x;
        this.y = y;
    }
}
\end{lstlisting}
\begin{lstlisting}[style=java, name=Alias.java]
public class Alias {
    public static void main(String[] args) {
        int n = 5;
        Punto p = new Punto(1, 2);
        Punto q = p;      // copia il riferimento, non l'oggetto
        q.x = 10;
        System.out.println(p.x + " " + q.x + " " + n);
    }
}
\end{lstlisting}
\begin{lstlisting}[style=shell, name=Terminale]
$ javac Punto.java Alias.java
$ java Alias
10 10 5
\end{lstlisting}
\end{minipage}\hfill
\begin{minipage}[c]{0.37\textwidth}
\centering
\begin{tikzpicture}
  \memblocco{swblue}{f}{(0,0)}{2.3cm}{main}{n/5, p/{}, q/{}}
  \memrif{f-2}
  \memrif{f-3}
  \memblocco{swteal}{pt}{(3.4,-0.25)}{1.9cm}{Punto}{x/10, y/2}
  \draw[memfreccia] (f-2-p) -- ([yshift=3pt]pt-sx);
  \draw[memfreccia] (f-3-p) -- ([yshift=-5pt]pt-sx);
  \memzona[swblue]{(f)(f-fondo)}{Stack}
  \memzona[swteal]{(pt)(pt-fondo)}{Heap}
\end{tikzpicture}

\vspace{4pt}
{\scriptsize La memoria alla fine di \code{main}.\par}
\end{minipage}

\medskip
\code{Punto q = p} non crea un secondo punto: copia nella variabile \code{q} lo stesso riferimento contenuto in \code{p}. Le
due variabili puntano allo stesso oggetto, quindi \code{q.x = 10} cambia anche \code{p.x}. Il valore di \code{n}, primitivo,
sta invece direttamente nel frame di \code{main}.
\end{esempio}

% ============================================================
\section{I pattern della Gang of Four}

I pattern della GoF sono divisi in tre famiglie, secondo il problema che risolvono:
\begin{itemize}
  \item \textbf{creazionali}: come si \textbf{creano} gli oggetti, nascondendo e controllando i dettagli della costruzione;
  \item \textbf{strutturali}: come si \textbf{compongono} classi e oggetti per formare strutture più grandi;
  \item \textbf{comportamentali}: come si distribuiscono \textbf{algoritmi e responsabilità} fra gli oggetti, e come
        questi interagiscono.
\end{itemize}

\begin{table}[H]
\centering
\renewcommand{\arraystretch}{1.2}
\newcommand{\studiato}[1]{\textcolor{swred}{\textbf{#1}}}
\begin{tabularx}{\textwidth}{P{3.0cm} L}
\toprule
\textbf{Pattern} & \textbf{Che cosa fa} \\
\midrule
\multicolumn{2}{@{}l}{\textit{Creazionali}} \\
\studiato{Builder} & Separa la costruzione di un oggetto complesso dalla sua rappresentazione. \\
Abstract Factory & Fornisce un'interfaccia per creare famiglie di oggetti correlati senza specificarne le classi concrete. \\
Factory Method & Definisce un'interfaccia per creare un oggetto, ma lascia alle sottoclassi la scelta della classe da
                 istanziare. \\
Prototype & Indica il tipo di oggetto da creare con un'istanza prototipo, e crea i nuovi oggetti copiando il prototipo. \\
\studiato{Singleton} & Garantisce che la classe abbia una sola istanza e fornisce un modo per accedervi. \\
\midrule
\multicolumn{2}{@{}l}{\textit{Strutturali}} \\
\studiato{Adapter} & Converte l'interfaccia di una classe in un'altra, così che due classi con interfacce diverse possano
                     lavorare insieme. \\
Bridge & Disaccoppia un'astrazione dalla sua implementazione, così che le due possano variare indipendentemente. \\
\studiato{Composite} & Compone gli oggetti in strutture ad albero per realizzare composizioni ricorsive. \\
\studiato{Decorator} & Aggiunge dinamicamente nuove responsabilità a un oggetto: è un'alternativa alle sottoclassi per
                       estenderne le funzionalità. \\
Facade & Fornisce un'interfaccia unificata a un sottosistema, per renderlo più facile da usare. \\
Proxy & Fornisce un surrogato di un oggetto per controllarne l'accesso. \\
Flyweight & Usa la condivisione per gestire in modo efficiente un gran numero di oggetti piccoli. \\
\midrule
\multicolumn{2}{@{}l}{\textit{Comportamentali}} \\
\studiato{Observer} & Definisce una dipendenza 1:N fra oggetti: quando uno cambia stato, gli altri vengono aggiornati
                      automaticamente. \\
\studiato{State} & Permette a un oggetto di cambiare comportamento a seconda del suo stato interno. \\
\studiato{Strategy} & Definisce una famiglia di algoritmi, li incapsula uno per uno e li rende intercambiabili, così che
                      possano variare indipendentemente da chi li usa. \\
Template Method & Definisce lo scheletro di un algoritmo, lasciando alle sottoclassi la definizione di alcuni passi. \\
Visitor & Rappresenta un'operazione da eseguire sugli elementi di una struttura di oggetti: si possono aggiungere
          operazioni senza modificare le classi degli elementi. \\
\bottomrule
\end{tabularx}
\caption{I pattern GoF presentati nelle slide. In rosso, come nelle slide, quelli che il corso approfondisce e che
vediamo in questo capitolo.}
\end{table}

Il catalogo originale ne conta 23: le slide riportano tutti i creazionali e tutti gli strutturali, e 5 degli 11
comportamentali (mancano, fra gli altri, \textit{Iterator}, quello degli iteratori delle collezioni Java, e
\textit{Command}).

% ============================================================
\section{Pattern creazionali}

In Java un oggetto si crea con \code{new} e un costruttore. È un meccanismo semplice ma rigido: il costruttore crea
\textbf{sempre} un oggetto nuovo, deve ricevere tutti i dati in una volta sola e lega chi lo chiama alla classe concreta. I
pattern creazionali \textbf{nascondono e controllano} la creazione degli oggetti; ne vediamo due, Singleton e Builder.

\subsection{Singleton}

\dfn{Singleton (creazionale)}{
  Il pattern \term{Singleton} costruisce classi che ammettono \textbf{una e una sola istanza} e forniscono un
  \textbf{punto di accesso globale} a quell'istanza.
}

Serve quando nel sistema deve esistere un solo oggetto di un certo tipo, e tutti devono poterlo raggiungere: la
connessione al database, la configurazione dell'applicazione, il registro dei log, la coda di stampa. Nel diagramma delle
classi i membri statici sono \underline{sottolineati}.

\Needspace{8\baselineskip}
\begin{center}
\begin{tikzpicture}
  \umlclass[4.2cm]{sg}{(0,0)}{Singleton}{\underline{- instance : Singleton}}{- Singleton()\\\underline{+ getInstance() : Singleton}}
  \node[notauml, anchor=west, text width=5.6cm] (sgn) at (3.0,0)
    {Il costruttore è privato: l'unico modo per avere l'oggetto è \texttt{getInstance()}, che lo crea alla prima chiamata
     e poi restituisce sempre lo stesso.};
  \draw[dashed, line width=0.5pt, draw=black!70] (sg.east) -- (sgn.west);
\end{tikzpicture}
\end{center}

\Needspace{24\baselineskip}
Ecco l'implementazione delle slide.

\begin{lstlisting}[style=java, name=SingletonClass.java]
public class SingletonClass {
    private static SingletonClass singletonObject = null;   // l'unica istanza

    private SingletonClass() {
        // codice di inizializzazione (facoltativo)
    }

    public static SingletonClass getSingletonObject() {
        if (singletonObject == null) {
            singletonObject = new SingletonClass();
        }
        return singletonObject;
    }

    @Override
    public Object clone() throws CloneNotSupportedException {
        throw new CloneNotSupportedException();
    }
}
\end{lstlisting}

\Needspace{12\baselineskip}
\begin{itemize}
  \item Il riferimento \code{singletonObject} è \textbf{static}: appartiene alla classe, non a un oggetto, quindi è uno solo
        e sopravvive mentre gli oggetti nascono e muoiono.
  \item È \textbf{private}: non può essere cambiato dall'esterno della classe (nessuno può scriverci \code{null} o un altro
        oggetto).
  \item Il \textbf{costruttore è privato}: lo può invocare solo il codice della classe stessa, cioè
        \code{getSingletonObject()}.
  \item Chiunque può ottenere il riferimento all'istanza chiamando il metodo \textbf{statico e pubblico}
        \code{getSingletonObject()}, che crea l'oggetto se non esiste ancora (\textit{lazy initialization},
        inizializzazione pigra).
  \item La possibilità di \textbf{copiare} l'oggetto è disabilitata: \code{clone()} è ridefinito in modo da lanciare sempre
        un'eccezione.
\end{itemize}

\begin{esempio}{Quanti oggetti esistono davvero?}
Il programma chiede l'istanza due volte e prova a copiarla.
\begin{lstlisting}[style=java, name=ProvaSingleton.java]
public class ProvaSingleton {
    public static void main(String[] args) {
        SingletonClass s1 = SingletonClass.getSingletonObject();
        SingletonClass s2 = SingletonClass.getSingletonObject();
        System.out.println("s1 == s2? " + (s1 == s2));   // stesso oggetto?

        try {
            Object copia = s1.clone();
        } catch (CloneNotSupportedException e) {
            System.out.println("Copia rifiutata: " + e);
        }
    }
}
\end{lstlisting}
\noindent\begin{minipage}[c]{0.48\textwidth}
\begin{lstlisting}[style=shell, name=Terminale]
$ javac SingletonClass.java ProvaSingleton.java
$ java ProvaSingleton
s1 == s2? true
Copia rifiutata: java.lang.CloneNotSupportedException
\end{lstlisting}
\end{minipage}\hfill
\begin{minipage}[c]{0.48\textwidth}
\centering
\begin{tikzpicture}
  \memblocco{swblue}{f}{(0,0)}{2.0cm}{main}{s1/{}, s2/{}}
  \memrif{f-1}
  \memrif{f-2}
  \memblocco{swteal}{o}{(3.1,-0.25)}{3.3cm}{SingletonClass}{}
  \memblocco{ph3}{st}{(3.1,-2.0)}{3.3cm}{SingletonClass}{singletonObject/{}}
  \memrif{st-1}
  \draw[memfreccia] (f-1-p) -- ([yshift=2pt]o.west);
  \draw[memfreccia] (f-2-p) -- ([yshift=-3pt]o.south west);
  \draw[memfreccia] (st-1-p) -- ++(0.6,0) |- (o.east);
  \memzona[swblue]{(f)(f-fondo)}{Stack}
  \memzona[swteal]{(o)(o-fondo)}{Heap}
  \memzona[ph3]{(st)(st-fondo)}{Campi statici}
\end{tikzpicture}
\end{minipage}

\medskip
La prima chiamata trova \code{singletonObject} a \code{null} e crea l'oggetto; la seconda lo trova già impostato e restituisce
lo stesso riferimento. Nel disegno ci sono tre riferimenti ma \textbf{un solo oggetto}, e \code{s1 == s2}, che confronta i
riferimenti, vale \code{true}.
\end{esempio}

Se un'altra classe prova a creare un secondo esemplare con \code{new}, a fermarla è già il \textbf{compilatore}:

\Needspace{16\baselineskip}
\begin{lstlisting}[style=java, name=Intruso.java]
public class Intruso {
    public static void main(String[] args) {
        SingletonClass altro = new SingletonClass();
    }
}
\end{lstlisting}
\begin{lstlisting}[style=shell, name=Terminale]
$ javac Intruso.java
Intruso.java:3: error: SingletonClass() has private access in SingletonClass
        SingletonClass altro = new SingletonClass();
                               ^
1 error
\end{lstlisting}

\Needspace{8\baselineskip}
\subsubsection*{L'esempio del corso: un nome che non cambia}

L'esempio distribuito con il corso (\code{EsempioSingleton}) chiama il metodo \code{getInstance} e gli passa un parametro,
che diventa una variabile di stato, \code{nome}. Il programma chiede l'istanza due volte, con due nomi diversi.

\Needspace{28\baselineskip}
\begin{lstlisting}[style=java, name=SingletonClass.java]
package singleton;

public class SingletonClass {
    private static SingletonClass instance = null;   // all'inizio non ci sono istanze
    private static String nome;                      // variabile di stato

    private SingletonClass(String nome) {            // costruttore privato
        setNome(nome);
    }

    public static SingletonClass getInstance(String nome) {
        if (instance == null) {
            instance = new SingletonClass(nome);
        }
        return instance;
    }

    public String getNome() {
        return nome;
    }

    private static void setNome(String nome) {
        SingletonClass.nome = nome;
    }
}
\end{lstlisting}

\Needspace{16\baselineskip}
\noindent\begin{minipage}[t]{0.55\textwidth}
\begin{lstlisting}[style=java, name=Esecutore.java]
package singleton;

public class Esecutore {
    public static void main(String[] args) {
        SingletonClass sc1;
        SingletonClass sc2;

        sc1 = SingletonClass.getInstance("uno");
        sc2 = SingletonClass.getInstance("due");

        System.out.println(sc1.getNome()
                           + " --- " + sc2.getNome());
    }
}
\end{lstlisting}
\end{minipage}\hfill
\begin{minipage}[t]{0.42\textwidth}
\begin{lstlisting}[style=shell, name=Terminale]
$ javac singleton/*.java
$ java singleton.Esecutore
uno --- uno
\end{lstlisting}
\end{minipage}

\Needspace{12\baselineskip}
\noindent\begin{minipage}[c]{0.5\textwidth}
Stampa \code{uno --- uno}, non \code{uno --- due}. La seconda chiamata trova \code{instance} già impostato: non entra
nell'\code{if}, non chiama il costruttore, e il parametro \code{"due"} viene semplicemente \textbf{ignorato}.
\code{sc1} e \code{sc2} sono lo stesso oggetto, e il nome è quello dato alla prima chiamata.

Nel disegno \code{nome} sta fra i campi statici, perché è dichiarato \code{static}; l'oggetto nello heap non ha campi
propri. Con un'istanza sola le due scelte si equivalgono, ma di solito lo stato si mette in un campo d'istanza.
\end{minipage}\hfill
\begin{minipage}[c]{0.46\textwidth}
\centering
\begin{tikzpicture}
  \memblocco{swblue}{f}{(0,0)}{2.0cm}{main}{sc1/{}, sc2/{}}
  \memrif{f-1}
  \memrif{f-2}
  \memblocco{swteal}{o}{(3.1,-0.25)}{3.0cm}{SingletonClass}{}
  \memblocco{ph3}{st}{(3.1,-2.0)}{3.0cm}{SingletonClass}{instance/{}, nome/"uno"}
  \memrif{st-1}
  \draw[memfreccia] (f-1-p) -- ([yshift=2pt]o.west);
  \draw[memfreccia] (f-2-p) -- ([yshift=-3pt]o.south west);
  \draw[memfreccia] (st-1-p) -- ++(0.6,0) |- (o.east);
  \memzona[swblue]{(f)(f-fondo)}{Stack}
  \memzona[swteal]{(o)(o-fondo)}{Heap}
  \memzona[ph3]{(st)(st-fondo)}{Campi statici}
\end{tikzpicture}
\end{minipage}

\warning{Un parametro che a volte non serve a nulla}{
  Un \code{getInstance} con parametri è ingannevole: chi scrive \code{getInstance("due")} si aspetta un oggetto di nome
  \code{"due"} e riceve, senza alcun avviso, quello creato da qualcun altro con un altro nome. Se il Singleton va
  configurato, è più chiaro farlo una volta sola, all'avvio, con un metodo dedicato (che segnala un errore se la
  configurazione è già stata fatta), e lasciare \code{getInstance()} senza parametri.
}

\Needspace{10\baselineskip}
\subsubsection*{Le conseguenze}

\begin{itemize}
  \item Della classe si può istanziare \textbf{al massimo un oggetto}. In alcune varianti si ammette invece un numero
        \textbf{fisso} di istanze (per esempio al massimo 3), tenute in un array di riferimenti che \code{getInstance}
        gestisce di conseguenza.
  \item Il Singleton è \textbf{migliore di un oggetto statico} (cioè di una classe con soli campi e metodi \code{static}):
        \begin{itemize}
          \item la creazione dell'oggetto si può \textbf{personalizzare}, nel costruttore e in \code{getInstance};
          \item la \textbf{copia} si può disabilitare;
          \item \textbf{non occupa memoria} finché l'oggetto non viene chiesto la prima volta.
        \end{itemize}
  \item È la soluzione che nel corso abbiamo già usato per l'oggetto \code{DatabaseConnection}.
\end{itemize}

\begin{esempio}{Al massimo tre istanze}
Una variante con un numero fisso di oggetti: tre stampanti, assegnate a turno a chi le chiede.

\begin{lstlisting}[style=java, name=Stampante.java]
public class Stampante {
    private static final int MAX = 3;
    private static final Stampante[] istanze = new Stampante[MAX];   // al massimo 3 oggetti
    private static int prossima = 0;
    private final int numero;

    private Stampante(int numero) {
        this.numero = numero;
    }

    public static Stampante getStampante() {
        int i = prossima;
        prossima = (prossima + 1) % MAX;          // a turno: 0, 1, 2, 0, 1, ...
        if (istanze[i] == null) {
            istanze[i] = new Stampante(i + 1);    // creata solo alla prima richiesta
        }
        return istanze[i];
    }

    public int getNumero() {
        return numero;
    }
}
\end{lstlisting}
\noindent\begin{minipage}[t]{0.55\textwidth}
\begin{lstlisting}[style=java, name=ProvaStampanti.java]
public class ProvaStampanti {
    public static void main(String[] args) {
        for (int k = 1; k <= 5; k++) {
            Stampante s = Stampante.getStampante();
            System.out.println("Richiesta " + k
                    + " -> stampante " + s.getNumero());
        }
    }
}
\end{lstlisting}
\end{minipage}\hfill
\begin{minipage}[t]{0.42\textwidth}
\begin{lstlisting}[style=shell, name=Terminale]
$ javac Stampante.java ProvaStampanti.java
$ java ProvaStampanti
Richiesta 1 -> stampante 1
Richiesta 2 -> stampante 2
Richiesta 3 -> stampante 3
Richiesta 4 -> stampante 1
Richiesta 5 -> stampante 2
\end{lstlisting}
\end{minipage}
\end{esempio}

\begin{esempio}{La connessione al database}
Una sola connessione condivisa da tutta l'applicazione, aperta solo quando serve la prima volta. È uno schema possibile
per \code{DatabaseConnection}; indirizzo e credenziali sono di esempio.
\begin{lstlisting}[style=java, name=DatabaseConnection.java]
import java.sql.Connection;
import java.sql.DriverManager;
import java.sql.SQLException;

public class DatabaseConnection {
    private static DatabaseConnection istanza = null;
    private final Connection connessione;

    private DatabaseConnection() throws SQLException {
        connessione = DriverManager.getConnection(
                "jdbc:postgresql://localhost:5432/negozio", "app", "segreta");
    }

    public static DatabaseConnection getInstance() throws SQLException {
        if (istanza == null) {
            istanza = new DatabaseConnection();      // la connessione si apre qui, una volta sola
        }
        return istanza;
    }

    public Connection getConnection() {
        return connessione;
    }
}
\end{lstlisting}
Ovunque serva il database si scrive \code{DatabaseConnection.getInstance().getConnection()}, senza passare la connessione
di classe in classe.
\end{esempio}

\info{{Singleton e thread}}{
  L'inizializzazione pigra delle slide non è sicura se più \textit{thread} chiamano \code{getSingletonObject()} nello
  stesso momento: entrambi possono trovare il riferimento a \code{null} ed entrambi creano un oggetto. Due rimedi comuni
  sono dichiarare il metodo \code{synchronized}, così che lo esegua un thread alla volta, oppure creare l'istanza subito,
  quando la classe viene caricata (cosa che la JVM fa una volta sola):
  \code{private static final SingletonClass ISTANZA = new SingletonClass();}. In quel caso si perde l'inizializzazione
  pigra.
}

\subsection{Builder}

Il Builder appartiene a un'intera famiglia di pattern creazionali (i \textit{Builder} e le \textit{Factory}) che affrontano
lo stesso problema: le classi \textbf{complesse}, per le quali il meccanismo di base dei costruttori non garantisce
abbastanza flessibilità ed efficienza.

\dfn{Builder (creazionale)}{
  Il pattern \term{Builder} rende la costruzione di un oggetto \textbf{flessibile e scomponibile} in passi, separandola dal
  costruttore vero e proprio dell'oggetto.
}

\Needspace{14\baselineskip}
\subsubsection*{Il problema: costruttori con troppi parametri}

Un computer ha molti componenti: alcuni obbligatori (processore, RAM), molti facoltativi (scheda video, disco, sistema
operativo), e in futuro se ne aggiungeranno altri. Con un solo costruttore, ogni creazione diventa un elenco di valori
senza nome:
\begin{lstlisting}[style=java]
Computer pc = new Computer("AMD Ryzen 7", 32, "AMD Radeon RX 7800", "SSD", 2000, true);
\end{lstlisting}
Che cosa sono \code{2000} e \code{true}? E se un componente manca, si passa \code{null}? Le alternative non sono migliori:
un costruttore per ogni combinazione di parametri facoltativi (costruttori che si moltiplicano a ogni nuova opzione), oppure
un costruttore vuoto seguito da tanti metodi \code{set}, che lascia esistere per un po' un computer montato a metà e lo
rende modificabile per sempre.

\Needspace{16\baselineskip}
\subsubsection*{La soluzione}

\begin{itemize}
  \item Il singolo oggetto \textbf{Product} è costruito dal \textbf{ConcreteBuilder}, un passo alla volta
        (\code{buildPart()}), e consegnato alla fine (\code{getResult()}).
  \item La classe \textbf{Director} può gestire la costruzione di un insieme di oggetti \code{Product}: conosce la sequenza
        dei passi e la ripete chiamando il Builder.
\end{itemize}

\begin{center}
\begin{tikzpicture}
  \umlclass[2.6cm]{bdir}{(0,0)}{Director}{}{+ construct()}
  \umlclass[2.6cm]{bbld}{(6.0,0)}{\textit{Builder}}{}{+ buildPart()}
  \umlclass[2.6cm]{bcon}{(6.0,-2.5)}{ConcreteBuilder}{}{+ buildPart()\\+ getResult()}
  \umlclass[2.2cm]{bpro}{(10.2,-2.5)}{Product}{}{}
  \draw[aggr] (bdir.east) -- node[molt, above]{builder} (bbld.west);
  \draw[gen] (bcon.north) -- (bbld.south);
  \draw[dep] (bcon.east) -- node[stereo, above]{«create»} (bpro.west);
  \node[notauml, text width=3.6cm, anchor=north] (bnota) at (0,-1.6) {for all objects in structure \{\\
    \hspace*{1em}builder.buildPart()\\\}};
  \draw[dashed, line width=0.5pt, draw=black!70] (bdir.south) -- (bnota.north);
\end{tikzpicture}
\end{center}

\Needspace{12\baselineskip}
Il Builder si usa quando:
\begin{itemize}
  \item la creazione di un prodotto complesso deve essere \textbf{indipendente dalle parti} che lo compongono; in
        particolare, il processo di creazione non deve conoscere il processo di \textbf{assemblaggio}, cioè come le parti
        vengono messe insieme per formare il prodotto;
  \item il processo di creazione deve permettere \textbf{rappresentazioni diverse} dell'oggetto costruito.
\end{itemize}
Gli esempi delle slide sono edifici: una casa con un piano, 3 stanze, 2 corridoi, 1 garage e tre porte; un grattacielo con
50 piani e, su ogni piano, 15 uffici e 5 corridoi. Gli stessi passi (\textit{aggiungi un piano}, \textit{aggiungi una
stanza}\ldots) producono edifici diversissimi.

\Needspace{12\baselineskip}
\subsubsection*{Un esempio: costruire un computer}

Vogliamo costruire oggetti della classe \code{Computer} per gestirne la pubblicità, la vendita e così via. Un computer ha
molti parametri, e in futuro potrebbe servirne qualcuno in più. Separiamo allora l'oggetto \code{Computer} da un oggetto
\code{ComputerBuilder}, che si occupa soltanto di gestire, in modo flessibile, tutte le varianti possibili della costruzione.
L'esempio, distribuito con il corso (\code{EsempioBuilder}), usa un ConcreteBuilder ma non il Director, che aggiungiamo più
avanti.

\begin{figure}[H]
\centering
\begin{tikzpicture}
  \umlclass[3.1cm]{cmain}{(-6.7,0)}{Main}{}{\underline{+ main(args : String[])}}
  \umlclass[6.0cm]{cbld}{(0,0)}{ComputerBuilder}
    {- cpu : String\\- ramGB : int\\- gpu : String = null\\- storageType : String = "Sconosciuto"\\- storageCapacityGB : int = 0\\- hasOS : boolean = false}
    {+ ComputerBuilder(cpu, ramGB)\\+ withDedicatedGPU(gpu) : ComputerBuilder\\+ withSSD(capacity) : ComputerBuilder\\+ withHDD(capacity) : ComputerBuilder\\+ includeOS() : ComputerBuilder\\+ build() : Computer}
  \umlclass[3.6cm]{cpc}{(6.9,0)}{Computer}
    {- cpu : String\\- ramGB : int\\- gpu : String\\- storageType : String\\- storageCapacityGB : int\\- hasOS : boolean}
    {\textasciitilde\ Computer(\ldots)\\+ toString() : String}
  \draw[dep] (cmain.east) -- node[stereo, above]{«use»} (cbld.west);
  \draw[dep] (cbld.east) -- node[stereo, above]{«create»} (cpc.west);
  \node[font=\tiny\itshape, text=black!60, anchor=south] at (cmain.north) {Client};
  \node[font=\tiny\itshape, text=black!60, anchor=south] at (cbld.north) {Builder concreto};
  \node[font=\tiny\itshape, text=black!60, anchor=south] at (cpc.north) {Prodotto complesso};
\end{tikzpicture}
\caption{Il diagramma delle classi dell'esempio. Il simbolo \textasciitilde\ indica la visibilità di package.}
\end{figure}

\Needspace{38\baselineskip}
Il \textbf{prodotto} ha tutti i campi \code{final}: una volta costruito, un computer non cambia più. Il suo costruttore non ha
modificatore di accesso (visibilità di \textit{package}): lo può chiamare solo il codice dello stesso package, cioè il
builder.

\begin{lstlisting}[style=java, name=Computer.java]
public class Computer {
    // Componenti obbligatori
    private final String cpu;
    private final int ramGB;

    // Componenti opzionali
    private final String gpu;
    private final String storageType;      // SSD o HDD
    private final int storageCapacityGB;
    private final boolean hasOS;           // sistema operativo

    // Costruttore package-private: lo usa solo ComputerBuilder (stesso package)
    Computer(String cpu, int ramGB, String gpu, String storageType,
             int storageCapacityGB, boolean hasOS) {
        this.cpu = cpu;
        this.ramGB = ramGB;
        this.gpu = gpu;
        this.storageType = storageType;
        this.storageCapacityGB = storageCapacityGB;
        this.hasOS = hasOS;
    }

    @Override
    public String toString() {
        return "Configurazione Computer:\n" +
               "  - CPU: " + cpu + "\n" +
               "  - RAM: " + ramGB + "GB\n" +
               "  - GPU: " + (gpu != null ? gpu : "Integrata") + "\n" +
               "  - Storage: " + (storageCapacityGB > 0
                       ? storageCapacityGB + "GB " + storageType : "Nessuno Storage") + "\n" +
               "  - OS Installato: " + hasOS;
    }
}
\end{lstlisting}

\Needspace{52\baselineskip}
Il \textbf{builder} riceve nel costruttore i parametri obbligatori e tiene quelli facoltativi, con i loro valori di
default, in campi modificabili. Ogni metodo \code{with\ldots} imposta un'opzione e restituisce \code{this}, cioè il builder
stesso: per questo le chiamate si possono \textbf{concatenare} (\textit{fluent interface}). Il metodo \code{build()}
controlla i dati e solo alla fine crea il prodotto.

\begin{lstlisting}[style=java, name=ComputerBuilder.java]
public class ComputerBuilder {
    // Parametri obbligatori: final, fissati dal costruttore del builder
    private final String cpu;
    private final int ramGB;

    // Parametri opzionali, con i valori di default
    private String gpu = null;
    private String storageType = "Sconosciuto";
    private int storageCapacityGB = 0;
    private boolean hasOS = false;

    public ComputerBuilder(String cpu, int ramGB) {
        this.cpu = cpu;
        this.ramGB = ramGB;
    }

    // Metodi "with": impostano un'opzione e restituiscono il builder stesso
    public ComputerBuilder withDedicatedGPU(String gpu) {
        this.gpu = gpu;
        return this;
    }

    public ComputerBuilder withSSD(int capacity) {
        this.storageType = "SSD";
        this.storageCapacityGB = capacity;
        return this;
    }

    public ComputerBuilder withHDD(int capacity) {
        this.storageType = "HDD";
        this.storageCapacityGB = capacity;
        return this;
    }

    public ComputerBuilder includeOS() {
        this.hasOS = true;
        return this;
    }

    // Metodo finale: controlla i dati e crea il Computer
    public Computer build() {
        if (ramGB < 4) {
            throw new IllegalStateException("RAM insufficiente. Minimo richiesto: 4GB.");
        }
        return new Computer(cpu, ramGB, gpu, storageType, storageCapacityGB, hasOS);
    }
}
\end{lstlisting}

\Needspace{28\baselineskip}
Il \textbf{client} ordina tre computer. Ogni ordine si legge come una frase: «un PC con un Ryzen 7 e 32 GB, con una GPU
dedicata, un SSD da 2000 GB e il sistema operativo».

\begin{lstlisting}[style=java, name=Main.java]
public class Main {
    public static void main(String[] args) {
        System.out.println("--- Ordine 1: PC Gaming ---");
        Computer gamingPC = new ComputerBuilder("AMD Ryzen 7", 32)
            .withDedicatedGPU("AMD Radeon RX 7800")
            .withSSD(2000)
            .includeOS()
            .build();
        System.out.println(gamingPC);

        System.out.println("\n--- Ordine 2: PC Ufficio Base ---");
        Computer officePC = new ComputerBuilder("Intel Core i3", 8)
            .withHDD(500)
            .build();                       // nessuna GPU dedicata, niente OS
        System.out.println(officePC);

        System.out.println("\n--- Ordine 3: PC obsoleto ---");
        Computer obsoletePC = new ComputerBuilder("Intel Core i3", 2)
            .withHDD(250)
            .build();
        System.out.println(obsoletePC);
    }
}
\end{lstlisting}

\Needspace{26\baselineskip}
\begin{lstlisting}[style=shell, name=Terminale]
$ javac *.java
$ java Main
--- Ordine 1: PC Gaming ---
Configurazione Computer:
  - CPU: AMD Ryzen 7
  - RAM: 32GB
  - GPU: AMD Radeon RX 7800
  - Storage: 2000GB SSD
  - OS Installato: true

--- Ordine 2: PC Ufficio Base ---
Configurazione Computer:
  - CPU: Intel Core i3
  - RAM: 8GB
  - GPU: Integrata
  - Storage: 500GB HDD
  - OS Installato: false

--- Ordine 3: PC obsoleto ---
Exception in thread "main" java.lang.IllegalStateException: RAM insufficiente. Minimo richiesto: 4GB.
        at ComputerBuilder.build(ComputerBuilder.java:43)
        at Main.main(Main.java:20)
\end{lstlisting}

Il terzo PC \textbf{non viene costruito}: \code{build()} lancia l'eccezione alla riga 43 di \code{ComputerBuilder.java}, e
il programma si ferma. La regola sui computer obsoleti non sta nella classe \code{Computer}, ma nel builder: il prodotto
resta semplice, e la logica di costruzione e di validazione è tutta in un posto.

\Needspace{30\baselineskip}
\subsubsection*{Che cosa succede in memoria}

\code{new ComputerBuilder(\ldots)} crea un oggetto nello heap. Ogni \code{with\ldots} lo modifica e ne restituisce il
riferimento (\code{return this}): la chiamata successiva della catena lavora quindi \textbf{sullo stesso oggetto}.
\code{build()} copia i valori in un oggetto \code{Computer} nuovo, ed è il suo riferimento che finisce in \code{gamingPC}.
Del builder, a quel punto, non resta alcun riferimento: è irraggiungibile, e il garbage collector lo eliminerà. Il
\code{Computer}, con i campi \code{final}, non si può più modificare.

\begin{center}
\begin{tikzpicture}
  \memblocco{swblue}{f}{(0,0)}{2.2cm}{main}{gamingPC/{}}
  \memrif{f-1}
  \memblocco{swteal}{c}{(3.1,0)}{4.3cm}{Computer}{cpu/"AMD Ryzen 7", ramGB/32, gpu/"AMD Radeon RX 7800",
    storageType/"SSD", storageCapacityGB/2000, hasOS/true}
  \memblocco[dashed]{swgray}{b}{(3.1,-4.35)}{4.3cm}{ComputerBuilder}{cpu/"AMD Ryzen 7", ramGB/32,
    gpu/"AMD Radeon RX 7800", storageType/"SSD", storageCapacityGB/2000, hasOS/true}
  \draw[memfreccia] (f-1-p) -- ([yshift=4pt]c.west);
  \draw[dep] (b.east) -- ++(0.45,0) |- node[stereo, pos=0.25, right]{build()} (c.east);
  \node[memnota, anchor=north west] at (0,-1.4) {il builder non è più\\raggiungibile da nessuna\\variabile: verrà
    eliminato};
  \memzona[swblue]{(f)(f-fondo)}{Stack}
  \memzona[swteal]{(c)(b-fondo)}{Heap}
\end{tikzpicture}
\end{center}

\Needspace{20\baselineskip}
\subsubsection*{Il Director: stessi passi, rappresentazioni diverse}

L'esempio del computer non usa il Director. Per vederlo all'opera servono un builder \textbf{astratto}, cioè
un'interfaccia con i passi di costruzione, e almeno due builder concreti. Il \code{Configuratore} (il Director) conosce le
«ricette» delle configurazioni ma non sa che cosa verrà prodotto: la stessa ricetta dà una scheda tecnica con un builder, un
preventivo con l'altro. Il metodo \code{getResult()} non sta nell'interfaccia, perché i due prodotti hanno tipi diversi
(\code{String} e \code{double}): è proprio come nel diagramma generale, dove \code{getResult()} compare solo nel
ConcreteBuilder.

\begin{center}
\begin{tikzpicture}
  \umlclass[3.8cm]{ddir}{(-0.4,0)}{Configuratore}{}{+ pcGaming(b : BuilderPC)\\+ pcUfficio(b : BuilderPC)}
  \umlclass[3.4cm]{dint}{(5.6,0)}{{\tiny\normalfont «interface»}\\BuilderPC}{}{+ cpu(modello, ramGB)\\+ gpu(modello)\\+ ssd(capacitaGB)}
  \umlclass[3.4cm]{dsch}{(3.6,-2.9)}{SchedaTecnica}{- testo : StringBuilder}{+ getResult() : String}
  \umlclass[3.0cm]{dpre}{(7.8,-2.9)}{Preventivo}{- totale : double}{+ getResult() : double}
  \draw[dep] (ddir.east) -- node[stereo, above]{«use»} (dint.west);
  \draw[realizza] (dsch.north) -- ++(0,0.45) -| ([xshift=-0.4cm]dint.south);
  \draw[realizza] (dpre.north) -- ++(0,0.45) -| ([xshift=0.4cm]dint.south);
\end{tikzpicture}
\end{center}

\Needspace{18\baselineskip}
\noindent\begin{minipage}[t]{0.49\textwidth}
\begin{lstlisting}[style=java, name=BuilderPC.java]
// Builder astratto: solo i passi di costruzione
public interface BuilderPC {
    void cpu(String modello, int ramGB);
    void gpu(String modello);
    void ssd(int capacitaGB);
}
\end{lstlisting}
\end{minipage}\hfill
\begin{minipage}[t]{0.49\textwidth}
\begin{lstlisting}[style=java, name=Configuratore.java]
// Director: conosce le ricette, non il prodotto
public class Configuratore {
    public void pcGaming(BuilderPC b) {
        b.cpu("AMD Ryzen 7", 32);
        b.gpu("AMD Radeon RX 7800");
        b.ssd(2000);
    }

    public void pcUfficio(BuilderPC b) {
        b.cpu("Intel Core i3", 8);
        b.ssd(256);
    }
}
\end{lstlisting}
\end{minipage}

\Needspace{22\baselineskip}
\noindent\begin{minipage}[t]{0.49\textwidth}
\begin{lstlisting}[style=java, name=SchedaTecnica.java]
// Primo ConcreteBuilder: un testo
public class SchedaTecnica implements BuilderPC {
    private final StringBuilder testo =
            new StringBuilder();

    public void cpu(String modello, int ramGB) {
        testo.append("CPU " + modello
                + ", RAM " + ramGB + " GB\n");
    }
    public void gpu(String modello) {
        testo.append("GPU " + modello + "\n");
    }
    public void ssd(int capacitaGB) {
        testo.append("SSD " + capacitaGB + " GB\n");
    }
    public String getResult() {
        return testo.toString();
    }
}
\end{lstlisting}
\end{minipage}\hfill
\begin{minipage}[t]{0.49\textwidth}
\begin{lstlisting}[style=java, name=Preventivo.java]
// Secondo ConcreteBuilder: un prezzo
public class Preventivo implements BuilderPC {
    private double totale = 0;

    public void cpu(String modello, int ramGB) {
        totale += 250 + 4 * ramGB;
    }
    public void gpu(String modello) {
        totale += 550;
    }
    public void ssd(int capacitaGB) {
        totale += 0.08 * capacitaGB;
    }
    public double getResult() {
        return totale;
    }
}
\end{lstlisting}
\end{minipage}

\Needspace{16\baselineskip}
\noindent\begin{minipage}[t]{0.62\textwidth}
\begin{lstlisting}[style=java, name=Negozio.java]
public class Negozio {
    public static void main(String[] args) {
        Configuratore director = new Configuratore();

        SchedaTecnica scheda = new SchedaTecnica();
        director.pcGaming(scheda);         // stessa ricetta...
        System.out.print(scheda.getResult());

        Preventivo preventivo = new Preventivo();
        director.pcGaming(preventivo);     // ...prodotto diverso
        System.out.println("Prezzo: " + preventivo.getResult()
                + " euro");
    }
}
\end{lstlisting}
\end{minipage}\hfill
\begin{minipage}[t]{0.35\textwidth}
\begin{lstlisting}[style=shell, name=Terminale]
$ javac *.java
$ java Negozio
CPU AMD Ryzen 7, RAM 32 GB
GPU AMD Radeon RX 7800
SSD 2000 GB
Prezzo: 1088.0 euro
\end{lstlisting}
\end{minipage}

\medskip
I prezzi sono di fantasia: $250 + 4 \cdot 32 = 378$ per processore e RAM, 550 per la scheda video,
$0{,}08 \cdot 2000 = 160$ per il disco, in tutto 1088 euro. Per aggiungere una nuova rappresentazione (un elenco dei pezzi
da ordinare al magazzino, per esempio) basta un nuovo ConcreteBuilder: il Director non cambia.

\Needspace{12\baselineskip}
\subsubsection*{Il Builder nelle librerie di Java}

Il pattern è molto usato nelle classi della libreria standard di Java:
\begin{itemize}
  \item \code{java.lang.StringBuilder} costruisce stringhe un pezzo alla volta (le \code{String} sono immutabili: senza un
        builder ogni concatenazione creerebbe una stringa nuova);
  \item \code{java.util.Locale.Builder} costruisce oggetti \code{Locale} complessi, definendo lingua, paese, sistema di
        scrittura e varianti;
  \item \code{java.time.format.DateTimeFormatterBuilder} costruisce formattatori di date e orari.
\end{itemize}

\Needspace{32\baselineskip}
\begin{lstlisting}[style=java, name=BuilderDiLibreria.java]
import java.time.LocalDateTime;
import java.time.format.DateTimeFormatter;
import java.time.format.DateTimeFormatterBuilder;
import java.util.Locale;

public class BuilderDiLibreria {
    public static void main(String[] args) {
        StringBuilder sb = new StringBuilder("Design");
        sb.append(' ').append("Patterns").append(", GoF ").append(1995);
        System.out.println(sb.toString());

        Locale locale = new Locale.Builder().setLanguage("it").setRegion("IT").build();
        System.out.println(locale + " = " + locale.getDisplayName(locale));

        DateTimeFormatter formato = new DateTimeFormatterBuilder()
                .appendPattern("dd/MM/yyyy")
                .appendLiteral(" alle ")
                .appendPattern("HH:mm")
                .toFormatter();
        System.out.println(LocalDateTime.of(2026, 10, 25, 9, 30).format(formato));
    }
}
\end{lstlisting}
\begin{lstlisting}[style=shell, name=Terminale]
$ java BuilderDiLibreria.java
Design Patterns, GoF 1995
it_IT = italiano (Italia)
25/10/2026 alle 09:30
\end{lstlisting}

Un builder l'avevamo già usato nel Capitolo~14, senza chiamarlo per nome: \code{HttpRequest.newBuilder(\ldots)} seguito da
\code{.header(\ldots)}, \code{.GET()} e infine \code{.build()}. E nel Capitolo~10 avevamo osservato che catene come
\code{builder.append("a").append("b")} non violano la Legge di Demetra, perché ogni chiamata restituisce lo stesso oggetto.

% ============================================================
\section{Pattern strutturali}

I pattern strutturali si occupano della \textbf{composizione} di classi e oggetti per creare strutture più grandi. Alcuni
usano l'ereditarietà (una classe può ereditare le proprietà di più superclassi indipendenti: in Java, estendendo una classe
e implementando delle interfacce), altri la composizione di oggetti. Permettono per esempio a librerie nate in modo
indipendente di lavorare insieme. Ne vediamo tre: Adapter, Composite e Decorator.

\subsection{Adapter}

\dfn{Adapter (strutturale)}{
  Il pattern \term{Adapter} (sinonimo: \textit{Wrapper}) converte l'interfaccia di una classe in un'altra interfaccia,
  quella che il client si aspetta. Permette così di far cooperare classi che altrimenti avrebbero interfacce incompatibili.
}

\textbf{Motivazione}: si vuole usare una classe già esistente e collaudata, ma la sua interfaccia non corrisponde a quella
che serve. È la situazione dell'adattatore da viaggio: la spina italiana funziona benissimo, ma la presa americana ha
un'altra forma; non si cambia né la spina né la presa, si mette in mezzo un adattatore.

\Needspace{8\baselineskip}
I partecipanti sono quattro:
\begin{itemize}
  \item \textbf{Client}: la classe che vuole usare il servizio da adattare;
  \item \textbf{Target}: l'interfaccia che vorremmo dal servizio;
  \item \textbf{Adaptee}: la classe esistente che fornisce il servizio;
  \item \textbf{Adapter}: la classe che adatta l'interfaccia fornita dall'Adaptee all'interfaccia Target.
\end{itemize}

L'adattamento si può realizzare in due modi: con l'\textbf{ereditarietà} (\textit{adapter di classe}: l'Adapter eredita
sia da Target sia da Adaptee) oppure con la \textbf{composizione} (\textit{adapter di oggetto}: l'Adapter eredita da Target e
contiene un riferimento a un Adaptee, a cui delega).

\begin{figure}[H]
\centering
\begin{tikzpicture}
  % adapter di classe
  \node[font=\scriptsize\bfseries] at (2.4,1.25) {Adapter di classe (ereditarietà)};
  \node[classe1=1.3cm] (ac) at (0,0) {Client};
  \umlclass[2.2cm]{at}{(2.4,0)}{\textit{Target}}{}{+ request()}
  \umlclass[2.4cm]{aa}{(5.2,0)}{Adaptee}{}{+ specificRequest()}
  \umlclass[2.2cm]{ad}{(2.4,-2.3)}{Adapter}{}{+ request()}
  \draw[dep] (ac.east) -- (at.west);
  \draw[gen] (ad.north) -- (at.south);
  \draw[gen] (ad.east) -- ++(0.9,0) -| (aa.south);
  \node[notauml, anchor=west] (an) at (4.1,-2.9) {\{ specificRequest(); \}};
  \draw[dashed, line width=0.5pt, draw=black!70] ([yshift=-0.25cm]ad.east) -- (an.west);
  % adapter di oggetto
  \begin{scope}[xshift=8.9cm]
    \node[font=\scriptsize\bfseries] at (2.4,1.25) {Adapter di oggetto (composizione)};
    \node[classe1=1.3cm] (bc) at (0,0) {Client};
    \umlclass[2.2cm]{bt}{(2.4,0)}{\textit{Target}}{}{+ request()}
    \umlclass[2.4cm]{ba}{(5.6,0)}{Adaptee}{}{+ specificRequest()}
    \umlclass[2.2cm]{bd}{(2.4,-2.3)}{Adapter}{}{+ request()}
    \draw[dep] (bc.east) -- (bt.west);
    \draw[gen] (bd.north) -- (bt.south);
    \draw[navig] (bd.east) -- node[molt, above, sloped, pos=0.55]{adaptee} (ba.south);
    \node[notauml, anchor=west] (bn) at (4.2,-2.95) {adaptee.specificRequest()};
    \draw[dashed, line width=0.5pt, draw=black!70] ([yshift=-0.25cm]bd.east) -- (bn.west);
  \end{scope}
\end{tikzpicture}
\caption{Le due forme dell'Adapter. In Java una classe estende una sola classe: l'adapter di classe si scrive
\code{class Adapter extends Adaptee implements Target}.}
\end{figure}

\Needspace{10\baselineskip}
L'Adapter si usa:
\begin{enumerate}
  \item se si vuole usare una classe esistente (l'Adaptee) la cui interfaccia è incompatibile con quella richiesta;
  \item se si vuole creare una classe (riusabile) che dovrà collaborare con classi non prevedibili al momento della sua
        creazione;
  \item se si vuole riusare un insieme di sottoclassi esistenti, ma è scomodo adattarne l'interfaccia creando una
        sottoclasse per ciascuna: meglio creare un adapter per la classe padre. Un adapter di oggetto può contenere
        un'istanza di qualunque sottoclasse dell'Adaptee, e le adatta tutte in una volta.
\end{enumerate}

\Needspace{16\baselineskip}
\subsubsection*{Un esempio: l'indice di massa corporea}

\noindent\begin{minipage}[c]{0.66\textwidth}
Stiamo realizzando un software a oggetti per il mercato statunitense, e serve una classe \code{Health} (il Client) che
fornisca, fra le altre cose, l'\textbf{indice di massa corporea} (IMC) di una persona, con l'interfaccia qui a destra. Non
conosciamo l'algoritmo per calcolare l'indice, e decidiamo di riusare la classe \code{IMC} di un software scritto per il
mercato italiano. Negli Stati Uniti però il peso si misura in \textbf{libbre} e l'altezza in \textbf{pollici}: sappiamo che
1 libbra vale 0,453 chilogrammi e 1 pollice 0,0254 metri. Vogliamo liberare chi programma \code{Health} dal dover conoscere
l'interfaccia e l'implementazione della classe \code{IMC}.
\end{minipage}\hfill
\begin{minipage}[c]{0.31\textwidth}
\centering
\begin{tikzpicture}
  \umlclass[4.2cm]{hh}{(0,0)}{Health}{+ IMC : double}{+ evaluateIMC(weight, height)}
  \umlclass[4.2cm]{hi}{(0,-1.9)}{IMC}{}{+ calcolaIMC(pesoKg : double,\\\hspace*{1.2em}altezzaM : double) : double}
  \node[font=\tiny\itshape, text=black!60, anchor=south] at (hh.north) {quello che serve (Client)};
  \node[font=\tiny\itshape, text=black!60, anchor=south] at (hi.north) {quello che esiste};
\end{tikzpicture}
\end{minipage}

\medskip
L'indice è il peso in chilogrammi diviso per il quadrato dell'altezza in metri: $\mathrm{IMC} = \text{peso} /
\text{altezza}^2$. Ecco la soluzione con un adapter di oggetto, nella stessa forma delle slide; \code{Adaptee} è la classe
che avevamo già a disposizione, che a sua volta usa il calcolo originale della classe \code{IMC}.

\Needspace{40\baselineskip}
\noindent\begin{minipage}[t]{0.49\textwidth}
\begin{lstlisting}[style=java, name=Main.java]
public class Main {
    public static void main(String[] args) {
        Health health = new Health();
        // 165 libbre, 70 pollici
        health.evaluateIMC(165, 70);
        System.out.println(health.IMC);
    }
}
\end{lstlisting}
\begin{lstlisting}[style=java, name=Health.java]
// Classe da realizzare (misure americane)
public class Health {
    public double IMC;

    public double evaluateIMC(double weight,
                              double height) {
        Target adapter = new Adapter();
        IMC = adapter.evaluate(weight, height);
        return IMC;
    }
}
\end{lstlisting}
\begin{lstlisting}[style=java, name=Target.java]
// Interfaccia che richiediamo
public interface Target {
    double evaluate(double weight, double height);
}
\end{lstlisting}
\end{minipage}\hfill
\begin{minipage}[t]{0.49\textwidth}
\begin{lstlisting}[style=java, name=Adapter.java]
// Classe che adatta la soluzione esistente
public class Adapter implements Target {
    private Adaptee adaptee = new Adaptee();

    @Override
    public double evaluate(double weight,
                           double height) {
        // libbre -> kg, pollici -> metri
        return adaptee.calcolaIMC(weight * 0.453,
                                  height * 0.0254);
    }
}
\end{lstlisting}
\begin{lstlisting}[style=java, name=Adaptee.java]
// Classe che avevamo gia' a disposizione
public class Adaptee {
    public double calcolaIMC(double pesoKg,
                             double altezzaM) {
        IMC imc = new IMC();
        return imc.calcolaIMC(pesoKg, altezzaM);
    }
}
\end{lstlisting}
\begin{lstlisting}[style=java, name=IMC.java]
// Calcolo originale (misure italiane)
public class IMC {
    public double calcolaIMC(double pesoKg,
                             double altezzaM) {
        return pesoKg / (altezzaM * altezzaM);
    }
}
\end{lstlisting}
\begin{lstlisting}[style=shell, name=Terminale]
$ javac *.java
$ java Main
23.643873818359882
\end{lstlisting}
\end{minipage}

\medskip
\code{Health} conosce soltanto l'interfaccia \code{Target}: non sa che esistono \code{Adaptee} e \code{IMC}, né che il calcolo
originale lavora in chilogrammi e metri. Tutta la conoscenza della «traduzione» è nell'\code{Adapter}, in una riga.

\warning{{Nelle slide le conversioni sono al contrario}}{
  Nella slide della soluzione l'adapter \textbf{divide}: \code{adaptee.calcolaIMC(weight/0.453, height/0.0254)}. Ma se
  1 libbra vale 0,453 kg, per passare dalle libbre ai chilogrammi bisogna \textbf{moltiplicare}: 165 libbre sono
  $165 \cdot 0{,}453 = 74{,}745$ kg, mentre $165 / 0{,}453 \simeq 364$ (sarebbero chili trasformati in libbre). Con le
  divisioni il programma stampa \code{4.795756183268009E-5}, cioè 0,000048: un indice senza senso. Con le moltiplicazioni
  stampa 23,64, il valore corretto per una persona di 75 kg alta 1,78 m. Qui il codice è corretto.
}

\Needspace{26\baselineskip}
\subsubsection*{Che cosa succede in memoria}

Il disegno fotografa lo stack nel momento più profondo, mentre la classe \code{IMC} esegue il calcolo. Ogni chiamata ha il suo
frame, e si vede dove avviene l'adattamento: nel frame di \code{Adapter.evaluate} i valori sono ancora in libbre e pollici,
in quello di \code{Adaptee.calcolaIMC} sono già in chilogrammi e metri. (Nei metodi d'istanza ogni frame contiene anche il
riferimento \code{this}, che qui non disegniamo.)

\begin{center}
\begin{tikzpicture}
  \memblocco{swblue}{s5}{(0,0)}{3.7cm}{IMC.calcolaIMC}{pesoKg/74.745, altezzaM/1.778}
  \memblocco{swblue}{s4}{(0,-1.6)}{3.7cm}{Adaptee.calcolaIMC}{pesoKg/74.745, altezzaM/1.778, imc/{}}
  \memrif{s4-3}
  \memblocco{swblue}{s3}{(0,-3.66)}{3.7cm}{Adapter.evaluate}{weight/165.0, height/70.0}
  \memblocco{swblue}{s2}{(0,-5.26)}{3.7cm}{Health.evaluateIMC}{weight/165.0, height/70.0, adapter/{}}
  \memrif{s2-3}
  \memblocco{swblue}{s1}{(0,-7.32)}{3.7cm}{main}{health/{}}
  \memrif{s1-1}
  % oggetti
  \memblocco{swteal}{oi}{(5.4,-2.55)}{2.2cm}{IMC}{}
  \memblocco{swteal}{oad}{(5.4,-6.15)}{2.4cm}{Adapter}{adaptee/{}}
  \memrif{oad-1}
  \memblocco{swteal}{oae}{(8.9,-6.4)}{2.2cm}{Adaptee}{}
  \memblocco{swteal}{oh}{(5.4,-7.6)}{2.4cm}{Health}{IMC/0.0}
  \draw[memfreccia] (s4-3-p) -- (oi.west);
  \draw[memfreccia] (s2-3-p) -- (oad.west);
  \draw[memfreccia] (oad-1-p) -- ([yshift=-2pt]oae.west);
  \draw[memfreccia] (s1-1-p) -- (oh.west);
  % la conversione
  \draw[-{Stealth[length=2.2mm]}, line width=1pt, swred] ([xshift=-0.45cm]s3-sx) to[bend left=55]
    node[left, font=\scriptsize, text=swred, align=right, xshift=-2pt]{libbre $\times\,0{,}453$\\pollici $\times\,0{,}0254$}
    ([xshift=-0.45cm]s4-sx);
  \memzona[swblue]{(s5)(s1-fondo)}{Stack (in cima la chiamata in corso)}
  \memzona[swteal]{(oi)(oh-fondo)(oae-fondo)}{Heap}
\end{tikzpicture}
\end{center}

\code{Health.IMC} vale ancora \code{0.0}: il valore verrà scritto solo quando le chiamate termineranno e i frame, uno dopo
l'altro dall'alto, spariranno restituendo il risultato.

\Needspace{14\baselineskip}
\subsubsection*{L'adapter di classe}

La stessa conversione si può scrivere con l'ereditarietà: l'adapter \textbf{è} un \code{Adaptee}, e chiama il metodo
ereditato. Il risultato è lo stesso.

\begin{lstlisting}[style=java, name=AdapterDiClasse.java]
// Adapter di classe: eredita dall'Adaptee invece di contenerlo
public class AdapterDiClasse extends Adaptee implements Target {
    @Override
    public double evaluate(double weight, double height) {
        return calcolaIMC(weight * 0.453, height * 0.0254);   // metodo ereditato
    }
}
\end{lstlisting}
\noindent\begin{minipage}[t]{0.55\textwidth}
\begin{lstlisting}[style=java, name=ProvaDiClasse.java]
public class ProvaDiClasse {
    public static void main(String[] args) {
        Target t = new AdapterDiClasse();
        System.out.println(t.evaluate(165, 70));
    }
}
\end{lstlisting}
\end{minipage}\hfill
\begin{minipage}[t]{0.42\textwidth}
\begin{lstlisting}[style=shell, name=Terminale]
$ javac *.java
$ java ProvaDiClasse
23.643873818359882
\end{lstlisting}
\end{minipage}

\medskip
L'adapter di classe è più compatto, ma è legato a \textbf{una} classe Adaptee e la espone tutta (chi ha un
\code{AdapterDiClasse} può chiamarne anche \code{calcolaIMC} direttamente). L'adapter di oggetto può adattare anche le
sottoclassi dell'Adaptee e nasconde completamente l'oggetto adattato: per questo, in Java, è la forma più comune.

\Needspace{14\baselineskip}
\subsubsection*{L'uso tipico: servizi e microservizi}

L'Adapter interviene di continuo quando si lavora con servizi o microservizi:
\begin{itemize}
  \item il progetto del nostro software ci porta a definire un'\textbf{interfaccia} (il Target, che può essere proprio
        un'interfaccia Java, perché chiede solo un servizio);
  \item un microservizio esistente fornisce il servizio che ci serve, ma con la \textbf{sua} interfaccia;
  \item scriviamo una classe Adapter che interagisce con la classe che fornisce il servizio (l'Adaptee), converte i
        parametri da lui forniti in quelli che interessano a noi, eventualmente riorganizzandoli ed elaborandoli, e li
        fornisce attraverso l'interfaccia Target.
\end{itemize}
Perché non modificare direttamente il servizio? Perché non l'abbiamo scritto noi e ne abbiamo solo l'interfaccia, oppure
perché è già usato da altri, che si aspettano l'interfaccia di oggi.

\begin{esempio}{Un convertitore di valuta sopra un servizio esterno}
Il nostro progetto vuole convertire importi in euro con una sola chiamata; la libreria del servizio dei cambi, che non
possiamo modificare, offre invece il tasso fra due valute qualsiasi. I tassi sono di esempio.

\begin{lstlisting}[style=java, name=ConvertitoreValuta.java]
// Target: l'interfaccia decisa dal nostro progetto
public interface ConvertitoreValuta {
    double inEuro(double importo, String valuta);
}
\end{lstlisting}
\begin{lstlisting}[style=java, name=ExchangeRateClient.java]
// Adaptee: il client di un servizio esterno, che non possiamo modificare
public class ExchangeRateClient {
    // quante unita' di "to" vale una unita' di "from"
    public double getRate(String from, String to) {
        // in realta' chiamerebbe il microservizio: GET /rates?from=USD&to=EUR
        return switch (from + "->" + to) {
            case "USD->EUR" -> 0.92;
            case "GBP->EUR" -> 1.17;
            default -> throw new IllegalArgumentException("Cambio non disponibile: " + from + "->" + to);
        };
    }
}
\end{lstlisting}
\begin{lstlisting}[style=java, name=AdapterCambi.java]
// Adapter: traduce la nostra domanda nella lingua del servizio
public class AdapterCambi implements ConvertitoreValuta {
    private final ExchangeRateClient client = new ExchangeRateClient();

    @Override
    public double inEuro(double importo, String valuta) {
        return importo * client.getRate(valuta, "EUR");
    }
}
\end{lstlisting}
\noindent\begin{minipage}[t]{0.6\textwidth}
\begin{lstlisting}[style=java, name=Cassa.java]
public class Cassa {
    public static void main(String[] args) {
        // il codice conosce solo il Target
        ConvertitoreValuta conv = new AdapterCambi();
        System.out.println("100 USD = "
                + conv.inEuro(100, "USD") + " EUR");
        System.out.println("100 GBP = "
                + conv.inEuro(100, "GBP") + " EUR");
    }
}
\end{lstlisting}
\end{minipage}\hfill
\begin{minipage}[t]{0.37\textwidth}
\begin{lstlisting}[style=shell, name=Terminale]
$ javac *.java
$ java Cassa
100 USD = 92.0 EUR
100 GBP = 117.0 EUR
\end{lstlisting}
\end{minipage}

\medskip
Se domani il servizio dei cambi cambia (o lo sostituiamo con un altro), si riscrive solo \code{AdapterCambi}: il resto del
programma continua a parlare con \code{ConvertitoreValuta}.
\end{esempio}

\subsection{Composite}

\dfn{Composite (strutturale)}{
  Il pattern \term{Composite} compone gli oggetti in \textbf{strutture ad albero} che rappresentano gerarchie
  tutto-parte, e permette ai client di trattare \textbf{allo stesso modo} gli oggetti composti e gli oggetti singoli.
}

\Needspace{16\baselineskip}
\noindent\begin{minipage}[c]{0.5\textwidth}
I partecipanti sono:
\begin{itemize}
  \item \textbf{Component}: dichiara l'interfaccia degli oggetti della composizione e il loro comportamento di default;
        dichiara anche un'interfaccia per accedere ai componenti figli e gestirli;
  \item \textbf{Leaf} (foglia): implementa gli oggetti senza figli e il loro comportamento;
  \item \textbf{Composite}: definisce il comportamento dei componenti con figli, memorizza i figli e implementa le
        operazioni sui figli definite dall'interfaccia Component;
  \item \textbf{Client}: usa gli oggetti della composizione attraverso l'interfaccia Component.
\end{itemize}
\end{minipage}\hfill
\begin{minipage}[c]{0.47\textwidth}
\centering
\begin{tikzpicture}
  \node[classe1=1.3cm] (cc) at (-2.6,0.6) {Client};
  \umlclass[3.6cm]{ccomp}{(0.9,0.6)}{{\tiny\normalfont «abstract»}\\Component}{}{+ operation()\\+ add(c : Component)\\+ remove(c : Component)\\+ getChild(i : int)}
  \umlclass[1.9cm]{cleaf}{(-1.1,-2.5)}{Leaf}{}{+ operation()}
  \umlclass[3.5cm]{ccmp}{(2.5,-2.5)}{Composite}{}{+ operation()\\+ add(c : Component)\\+ remove(c : Component)\\+ getChild(i : int)}
  \draw[navig] (cc.east) -- (ccomp.west);
  \draw[gen] (cleaf.north) -- ++(0,0.55) -| ([xshift=-0.5cm]ccomp.south);
  \draw[gen] (ccmp.north) -- ++(0,0.35) -| ([xshift=0.5cm]ccomp.south);
  \draw[aggr] (ccmp.east) -- ++(0.45,0) |- node[molt, pos=0.75, above]{figli 0..*} (ccomp.east);
\end{tikzpicture}
\end{minipage}

\medskip
La relazione decisiva è l'aggregazione da \code{Composite} a \code{Component}: un composto contiene dei \textbf{componenti}, e
quindi può contenere sia foglie sia altri composti. È questo che rende la struttura ricorsiva, e permette alberi di
qualunque profondità.

\Needspace{14\baselineskip}
\subsubsection*{Un esempio: il file system}

Vogliamo realizzare un albero di directory e file, con una funzione \code{ls} che li elenca rispettando il
\textbf{rientro}: ogni livello dell'albero si sposta un po' più a destra. Una \code{Directory} può contenere file e altre
directory, eventualmente vuote; \code{AbstractFile} rappresenta tutto ciò che può stare in una directory. L'esempio è quello
distribuito con il corso (\code{EsempioComposite}).

\begin{figure}[H]
\centering
\begin{tikzpicture}
  \umlclass[4.1cm]{fe}{(0,0)}{EsempioComposite}{\underline{+ indentazione : StringBuffer}}{\underline{+ main(args : String[])}}
  \umlclass[4.8cm]{fd}{(6.0,0)}{Directory}{- nome : String\\- listaFile : ArrayList<AbstractFile>}{+ add(obj : Object)\\+ ls()}
  \umlclass[2.2cm]{ff}{(10.9,0)}{File}{- nome : String}{+ ls()}
  \umlclass[2.2cm]{fa}{(8.4,-2.6)}{{\tiny\normalfont «interface»}\\AbstractFile}{}{+ ls()}
  \draw[dep] (fe.east) -- node[stereo, above]{«create»} (fd.west);
  \draw[dep] (fe.north) -- ++(0,0.5) -| node[stereo, above, pos=0.25]{«create»} (ff.north);
  \draw[realizza] (fd.south) -- ++(0,-0.6) -| ([xshift=-0.4cm]fa.north);
  \draw[realizza] (ff.south) -- ++(0,-0.6) -| ([xshift=0.4cm]fa.north);
  \draw[compos] ([yshift=-0.35cm]fd.south west) |- node[molt, pos=0.75, below]{listaFile *} (fa.west);
\end{tikzpicture}
\caption{Il diagramma delle classi dell'esempio: \code{Directory} è il Composite, \code{File} la Leaf, \code{AbstractFile} il
Component.}
\end{figure}

\Needspace{16\baselineskip}
\noindent\begin{minipage}[t]{0.36\textwidth}
\begin{lstlisting}[style=java, name=AbstractFile.java]
interface AbstractFile {
    public void ls();
}
\end{lstlisting}
\end{minipage}\hfill
\begin{minipage}[t]{0.61\textwidth}
\begin{lstlisting}[style=java, name=File.java]
class File implements AbstractFile {
    private String nome;

    public File(String n) {
        nome = n;
    }

    public void ls() {
        System.out.println(EsempioComposite.indentazione + nome);
    }
}
\end{lstlisting}
\end{minipage}

\Needspace{30\baselineskip}
La \code{Directory} tiene i figli in una lista di \code{AbstractFile}. Il suo \code{ls()} stampa il proprio nome, aggiunge tre
spazi al rientro, chiama \code{ls()} su ogni figlio (senza sapere se è un file o una directory: è il
\textbf{polimorfismo}), e alla fine toglie i tre spazi.

\begin{lstlisting}[style=java, name=Directory.java]
import java.util.ArrayList;

class Directory implements AbstractFile {
    private String nome;
    private ArrayList<AbstractFile> listaFile = new ArrayList<AbstractFile>();

    public Directory(String name) {
        nome = name;
    }

    public void add(Object obj) {
        listaFile.add((AbstractFile) obj);
    }

    public void ls() {
        System.out.println(EsempioComposite.indentazione + nome);
        EsempioComposite.indentazione.append("   ");
        for (int i = 0; i < listaFile.size(); ++i) {
            AbstractFile obj = listaFile.get(i);
            obj.ls();                                  // polimorfismo
        }
        EsempioComposite.indentazione.setLength(EsempioComposite.indentazione.length() - 3);
    }
}
\end{lstlisting}

\Needspace{26\baselineskip}
\noindent\begin{minipage}[t]{0.7\textwidth}
\begin{lstlisting}[style=java, name=EsempioComposite.java]
public class EsempioComposite {
    public static StringBuffer indentazione = new StringBuffer();

    public static void main(String[] args) {
        Directory one = new Directory("dir111"),
                  two = new Directory("dir222"),
                  thr = new Directory("dir333");
        File a = new File("a"), b = new File("b"), c = new File("c"),
             d = new File("d"), e = new File("e");
        one.add(a);
        one.add(two);
        one.add(b);
        two.add(c);
        two.add(d);
        two.add(thr);
        thr.add(e);
        one.ls();
    }
}
\end{lstlisting}
\end{minipage}\hfill
\begin{minipage}[t]{0.27\textwidth}
\begin{lstlisting}[style=shell, name=Terminale]
$ javac *.java
$ java EsempioComposite
dir111
   a
   dir222
      c
      d
      dir333
         e
   b
\end{lstlisting}
\end{minipage}

\medskip
Il \code{main} chiama \code{ls()} \textbf{una sola volta}, sulla radice: l'albero intero viene visitato dalla ricorsione.
Rispetto alle slide, il codice del corso usa una lista generica (\code{ArrayList<AbstractFile>}), così che \code{get(i)}
restituisce già un \code{AbstractFile} senza conversioni.

\Needspace{30\baselineskip}
\subsubsection*{Che cosa succede in memoria}

A destra c'è l'albero degli oggetti costruito dalle \code{add}; a sinistra lo stack nel momento in cui viene stampata la riga
di \code{e}, la più profonda. Per non affollare il disegno, nei frame il riferimento \code{this} è indicato con il nome
dell'oggetto puntato, e le caselle della lista \code{listaFile} (che è un oggetto \code{ArrayList} a parte) sono disegnate
dentro la directory.

\begin{center}
\begin{tikzpicture}
  % lo stack: in cima la chiamata in corso
  \memblocco{swblue}{k5}{(0,0)}{2.9cm}{File.ls()}{this/e}
  \memblocco{swblue}{k4}{(0,-1.12)}{2.9cm}{Directory.ls()}{this/dir333, i/0}
  \memblocco{swblue}{k3}{(0,-2.7)}{2.9cm}{Directory.ls()}{this/dir222, i/2}
  \memblocco{swblue}{k2}{(0,-4.28)}{2.9cm}{Directory.ls()}{this/dir111, i/1}
  \memblocco{swblue}{k1}{(0,-5.86)}{2.9cm}{main}{one/dir111}
  % il campo statico
  \memblocco{ph3}{kst}{(0,-7.95)}{4.6cm}{EsempioComposite}{indentazione/{"\textvisiblespace\textvisiblespace\textvisiblespace\textvisiblespace\textvisiblespace\textvisiblespace\textvisiblespace\textvisiblespace\textvisiblespace"}}
  % l'albero nello heap
  \memblocco{swteal}{t1}{(4.3,0)}{2.5cm}{Directory}{nome/"dir111", {[0]}/{}, {[1]}/{}, {[2]}/{}}
  \memrif{t1-2}
  \memrif{t1-3}
  \memrif{t1-4}
  \memblocco{swteal}{ta}{(8.0,0.2)}{2.1cm}{File}{nome/"a"}
  \memblocco{swteal}{t2}{(8.0,-1.3)}{2.5cm}{Directory}{nome/"dir222", {[0]}/{}, {[1]}/{}, {[2]}/{}}
  \memrif{t2-2}
  \memrif{t2-3}
  \memrif{t2-4}
  \memblocco{swteal}{tb}{(8.0,-4.15)}{2.1cm}{File}{nome/"b"}
  \memblocco{swteal}{tc}{(11.7,-0.3)}{2.1cm}{File}{nome/"c"}
  \memblocco{swteal}{td}{(11.7,-1.65)}{2.1cm}{File}{nome/"d"}
  \memblocco{swteal}{t3}{(11.7,-3.0)}{2.5cm}{Directory}{nome/"dir333", {[0]}/{}}
  \memrif{t3-2}
  \memblocco{swteal}{te}{(15.0,-3.0)}{2.1cm}{File}{nome/"e"}
  \draw[memfreccia] (t1-2-p) -- ++(0.45,0) |- (ta.west);
  \draw[memfreccia] (t1-3-p) -- ++(0.45,0) |- (t2.west);
  \draw[memfreccia] (t1-4-p) -- ++(0.45,0) |- (tb.west);
  \draw[memfreccia] (t2-2-p) -- ++(0.45,0) |- (tc.west);
  \draw[memfreccia] (t2-3-p) -- ++(0.45,0) |- (td.west);
  \draw[memfreccia] (t2-4-p) -- ++(0.45,0) |- (t3.west);
  \draw[memfreccia] (t3-2-p) -- ++(0.45,0) |- (te.west);
  \memzona[swblue]{(k5)(k1-fondo)}{Stack}
  \memzona[ph3]{(kst)(kst-fondo)}{Campi statici}
  \memzona[swteal]{(t1)(ta)(te-fondo)(tc)(tb-fondo)}{Heap}
\end{tikzpicture}
\end{center}

Ogni chiamata di \code{ls()} su una directory ha il suo frame, con il \textbf{suo} indice \code{i}: \code{dir111} è fermo al
figlio 1 (\code{dir222}), \code{dir222} al figlio 2 (\code{dir333}), \code{dir333} al figlio 0 (\code{e}). Il rientro, condiviso
da tutti perché statico, in quel momento vale nove spazi: tre per ogni directory aperta. Quando \code{e.ls()} termina, il suo
frame sparisce e riprende il ciclo di \code{dir333}, che toglie tre spazi e termina a sua volta; poi riprende \code{dir222},
che ha finito i figli; infine \code{dir111} passa al figlio 2 e stampa \code{b} con tre spazi soltanto.

\info{{Trasparenza o sicurezza?}}{
  Nel diagramma generale le operazioni sui figli (\code{add}, \code{remove}, \code{getChild}) sono dichiarate nel
  Component: il client può trattare foglie e composti nello stesso identico modo (\textbf{trasparenza}), ma può anche
  chiedere a una foglia di aggiungere un figlio, cosa che non ha senso e va gestita, per esempio con un'eccezione.
  Nell'esempio del file system \code{add} esiste solo in \code{Directory}: il compilatore impedisce di aggiungere un figlio
  a un \code{File} (\textbf{sicurezza}), ma il client deve sapere quando ha in mano una directory. È un compromesso, e
  nessuna delle due scelte è sbagliata. Con \code{add(AbstractFile f)} invece di \code{add(Object obj)}, infine, il
  compilatore rifiuterebbe anche gli oggetti che non sono né file né directory, che oggi passano la compilazione e
  falliscono solo a tempo di esecuzione, sulla conversione \code{(AbstractFile) obj}.
}

\begin{esempio}{Il rientro come parametro, invece che come variabile statica}
La \code{StringBuffer} statica è uno \textbf{stato globale}: tutte le chiamate la modificano, e ogni \code{ls()} deve
ricordarsi di rimettere le cose a posto alla fine. Si può fare a meno, passando il rientro come parametro: ogni chiamata
riceve il rientro del suo livello, e non c'è nulla da ripristinare.

\begin{lstlisting}[style=java, name=Directory.java]
import java.util.ArrayList;
import java.util.List;

class Directory implements AbstractFile {
    private final String nome;
    private final List<AbstractFile> figli = new ArrayList<>();
    public Directory(String nome) { this.nome = nome; }

    public void add(AbstractFile f) {       // solo file e directory
        figli.add(f);
    }

    public void ls(String rientro) {
        System.out.println(rientro + nome);
        for (AbstractFile f : figli) {
            f.ls(rientro + "   ");          // tre spazi in piu'
        }
    }
}
\end{lstlisting}
\noindent\begin{minipage}[t]{0.49\textwidth}
\begin{lstlisting}[style=java, name=AbstractFile.java]
interface AbstractFile {
    void ls(String rientro);
}
\end{lstlisting}
\end{minipage}\hfill
\begin{minipage}[t]{0.49\textwidth}
\begin{lstlisting}[style=java, name=File.java]
class File implements AbstractFile {
    private final String nome;
    public File(String nome) {
        this.nome = nome;
    }

    public void ls(String rientro) {
        System.out.println(rientro + nome);
    }
}
\end{lstlisting}
\end{minipage}

\medskip
La radice si stampa con \code{one.ls("")}, e l'output è identico a quello di prima.
\end{esempio}

\Needspace{14\baselineskip}
\subsubsection*{Le conseguenze}

\begin{itemize}
  \item Il Composite permette di definire gerarchie di classi fatte di oggetti \textbf{primitivi} e \textbf{composti}: gli
        oggetti primitivi si compongono in oggetti più complessi, che a loro volta si compongono ricorsivamente. Ovunque il
        client si aspetta un oggetto primitivo, si può usare indifferentemente un oggetto composto.
  \item \textbf{Semplifica il client}: i client usano in modo uniforme le strutture composte e gli oggetti singoli, e di
        solito non sanno (né dovrebbero preoccuparsi di sapere) se stanno lavorando con una foglia o con un composto.
  \item È \textbf{flessibile}: rende facile aggiungere nuovi tipi di componenti; nuove sottoclassi Leaf o Composite
        funzionano automaticamente con le strutture e con il codice dei client esistenti.
  \item È spesso usato per la gerarchia dei \textbf{widget} di un'interfaccia grafica. In Java AWT e Swing, per esempio,
        un \code{Container} (una finestra, un pannello) è a sua volta un \code{Component} e contiene altri
        \code{Component}: pulsanti, etichette e altri pannelli.
\end{itemize}

\subsection{Decorator}

Il pattern Decorator permette di creare oggetti complessi \textbf{aggiungendo funzionalità} a oggetti più semplici.
\begin{itemize}
  \item L'oggetto complesso \textbf{include} l'oggetto semplice e lo «decora», aggiungendo funzionalità: l'oggetto originale
        viene passato come parametro al costruttore del decoratore, e si ottiene l'oggetto «decorato».
  \item Più decoratori si possono anche \textbf{concatenare}, aggiungendo funzionalità un passo alla volta: l'oggetto
        complesso è l'ultimo anello della catena.
  \item La concatenazione può avvenire secondo una composizione \textbf{arbitraria}: il numero dei comportamenti possibili
        cresce in modo combinatorio rispetto al numero dei decoratori disponibili.
\end{itemize}
È un'\textbf{alternativa all'ereditarietà}, singola o multipla. Con l'ereditarietà le funzionalità si aggiungono in modo
\textbf{statico}, secondo la gerarchia delle classi scritta nel codice: non si può ottenere a tempo di esecuzione una
combinazione arbitraria di funzionalità, né aggiungerne o toglierne.

\dfn{Decorator (strutturale)}{
  Il pattern \term{Decorator} aggiunge \textbf{dinamicamente} nuove responsabilità a un oggetto, avvolgendolo in un altro
  oggetto che ha la stessa interfaccia: è un'alternativa flessibile alle sottoclassi per estenderne le funzionalità.
}

\Needspace{16\baselineskip}
\noindent\begin{minipage}[c]{0.45\textwidth}
\centering
\begin{tikzpicture}
  \umlclass[2.4cm]{dco}{(1.4,0)}{\textit{Component}}{}{+ operation()}
  \umlclass[3.1cm]{dcc}{(-0.6,-2.3)}{ConcreteComponent}{}{+ operation()}
  \umlclass[2.6cm]{dde}{(3.2,-2.3)}{\textit{Decorator}}{- component}{+ operation()}
  \umlclass[3.1cm]{dcd}{(3.2,-4.8)}{ConcreteDecorator}{}{+ operation()}
  \draw[gen] (dcc.north) -- ++(0,0.5) -| ([xshift=-0.4cm]dco.south);
  \draw[gen] (dde.north) -- ++(0,0.5) -| ([xshift=0.4cm]dco.south);
  \draw[aggr] (dde.east) -- ++(0.5,0) |- (dco.east);
  \draw[gen] (dcd.north) -- (dde.south);
\end{tikzpicture}
\end{minipage}\hfill
\begin{minipage}[c]{0.52\textwidth}
\begin{itemize}
  \item \textbf{Component} è una classe astratta o un'interfaccia.
  \item \textbf{ConcreteComponent} è l'implementazione di base del Component.
  \item \textbf{Decorator} è una classe astratta che \textbf{include} un Component (lo ha come attributo) e ne
        \textbf{eredita} l'operazione: è un Component che contiene un Component.
  \item \textbf{ConcreteDecorator} migliora il Component: può assumerne tutte le caratteristiche, perché ha un attributo
        \code{component} a cui delegare, e può modificarne i metodi ridefinendoli (\textit{override}).
\end{itemize}
\end{minipage}

\Needspace{12\baselineskip}
\subsubsection*{Un esempio: la pizza}

Una pizza base (olio e pomodoro) costa 3 euro; ogni ingrediente aggiunto la rende più ricca e più cara. Invece di scrivere
una classe per ogni combinazione, scriviamo una classe per ogni \textbf{ingrediente}, che decora una pizza qualsiasi.
L'esempio è quello distribuito con il corso (\code{EsempioDecorator}).

\begin{figure}[H]
\centering
\begin{tikzpicture}
  \umlclass[3.7cm]{pz}{(0,0)}{{\tiny\normalfont «interface»}\\Pizza}{}{+ getDescription() : String\\+ getCost() : double}
  \umlclass[3.7cm]{pzb}{(-4.6,-2.6)}{OlioEPomodoro}{}{+ getDescription() : String\\+ getCost() : double}
  \umlclass[4.4cm]{pzd}{(3.2,-2.6)}{\textit{PizzaDecorator}}{\# pizza : Pizza}{+ PizzaDecorator(pizza : Pizza)\\+ getDescription() : String\\+ getCost() : double}
  \umlclass[2.6cm]{pzm}{(-0.2,-5.7)}{Mozzarella}{}{+ getDescription()\\+ getCost()}
  \umlclass[2.6cm]{pzp}{(3.2,-5.7)}{Prosciutto}{}{+ getDescription()\\+ getCost()}
  \umlclass[2.6cm]{pzf}{(6.6,-5.7)}{Fungo}{}{+ getDescription()\\+ getCost()}
  \node[classe1=1.4cm] (pzc) at (-4.6,0) {Main};
  \draw[realizza] (pzb.north) -- ++(0,0.6) -| ([xshift=-0.5cm]pz.south);
  \draw[realizza] (pzd.north) -- ++(0,0.6) -| ([xshift=0.5cm]pz.south);
  \draw[aggr] (pzd.east) -- ++(0.6,0) |- node[molt, pos=0.75, above]{decora} (pz.east);
  \foreach \x in {pzm, pzp, pzf}{ \draw[gen] (\x.north) -- ++(0,0.45) -| (pzd.south); }
  \draw[dep] (pzc.east) -- node[stereo, above]{«use»} (pz.west);
\end{tikzpicture}
\caption{Il diagramma delle classi della pizza: \code{Pizza} è il Component, \code{OlioEPomodoro} il ConcreteComponent,
\code{PizzaDecorator} il Decorator, gli ingredienti i ConcreteDecorator.}
\end{figure}

\Needspace{24\baselineskip}
\noindent\begin{minipage}[t]{0.49\textwidth}
\begin{lstlisting}[style=java, name=Pizza.java]
// Component: l'interfaccia comune
public interface Pizza {
    String getDescription();
    double getCost();
}
\end{lstlisting}
\begin{lstlisting}[style=java, name=OlioEPomodoro.java]
// ConcreteComponent: la pizza base
public class OlioEPomodoro implements Pizza {
    @Override
    public String getDescription() {
        return "Pizza Olio e Pomodoro";
    }

    @Override
    public double getCost() {
        return 3.00;
    }
}
\end{lstlisting}
\end{minipage}\hfill
\begin{minipage}[t]{0.49\textwidth}
\begin{lstlisting}[style=java, name=PizzaDecorator.java]
// Decorator: e' una Pizza e contiene una Pizza
public abstract class PizzaDecorator
        implements Pizza {
    protected Pizza pizza;       // la pizza avvolta

    public PizzaDecorator(Pizza pizza) {
        this.pizza = pizza;
    }

    @Override
    public String getDescription() {
        return pizza.getDescription();   // delega
    }

    @Override
    public double getCost() {
        return pizza.getCost();          // delega
    }
}
\end{lstlisting}
\end{minipage}

\Needspace{20\baselineskip}
Ogni ingrediente è un decoratore concreto: aggiunge il proprio nome alla descrizione e il proprio prezzo al costo della pizza
che avvolge. \code{Prosciutto} (1 euro) e \code{Fungo} (2 euro, descrizione «Funghi») sono identici a \code{Mozzarella},
tranne che per il testo e il prezzo.

\noindent\begin{minipage}[t]{0.49\textwidth}
\begin{lstlisting}[style=java, name=Mozzarella.java]
// ConcreteDecorator: aggiunge la mozzarella
public class Mozzarella extends PizzaDecorator {
    public Mozzarella(Pizza pizza) {
        super(pizza);
    }

    @Override
    public String getDescription() {
        return pizza.getDescription()
                + ", Mozzarella";
    }

    @Override
    public double getCost() {
        return pizza.getCost() + 1.50;
    }
}
\end{lstlisting}
\end{minipage}\hfill
\begin{minipage}[t]{0.49\textwidth}
\begin{lstlisting}[style=java, name=Fungo.java]
// ConcreteDecorator: aggiunge i funghi
public class Fungo extends PizzaDecorator {
    public Fungo(Pizza pizza) {
        super(pizza);
    }

    @Override
    public String getDescription() {
        return pizza.getDescription()
                + ", Funghi";
    }

    @Override
    public double getCost() {
        return pizza.getCost() + 2.00;
    }
}
\end{lstlisting}
\end{minipage}

\Needspace{34\baselineskip}
Il \code{main} ordina una pizza base, poi la stessa con la mozzarella, poi una pizza «personalizzata» con quattro ingredienti,
fra cui la mozzarella due volte.

\begin{lstlisting}[style=java, name=Main.java]
public class Main {
    public static void main(String[] args) {
        // Ordiniamo una pizza base
        Pizza myPizza = new OlioEPomodoro();
        System.out.println("Pizza ordinata: " + myPizza.getDescription());
        System.out.println("Costo: €" + myPizza.getCost());
        System.out.println("--------------------");

        // Aggiungiamo la mozzarella alla pizza
        myPizza = new Mozzarella(myPizza);
        System.out.println("Pizza ordinata: " + myPizza.getDescription());
        System.out.println("Costo: €" + myPizza.getCost());
        System.out.println("--------------------");

        // Una pizza piu' ricca, con doppia mozzarella
        Pizza customPizza = new OlioEPomodoro();
        customPizza = new Prosciutto(customPizza);
        customPizza = new Mozzarella(customPizza);
        customPizza = new Fungo(customPizza);
        customPizza = new Mozzarella(customPizza);
        System.out.println("\n*** La tua pizza personalizzata ***");
        System.out.println("Descrizione: " + customPizza.getDescription());
        System.out.println("Costo totale: €" + customPizza.getCost());
    }
}
\end{lstlisting}
\begin{lstlisting}[style=shell, name=Terminale]
$ javac *.java
$ java Main
Pizza ordinata: Pizza Olio e Pomodoro
Costo: €3.0
--------------------
Pizza ordinata: Pizza Olio e Pomodoro, Mozzarella
Costo: €4.5
--------------------

*** La tua pizza personalizzata ***
Descrizione: Pizza Olio e Pomodoro, Prosciutto, Mozzarella, Funghi, Mozzarella
Costo totale: €9.0
\end{lstlisting}

\Needspace{24\baselineskip}
\subsubsection*{Che cosa succede in memoria}

Ogni \code{new} aggiunge un anello alla catena: il nuovo decoratore riceve nel costruttore la pizza costruita fin lì e la
conserva nel campo \code{pizza}; il riferimento al nuovo anello rimpiazza quello vecchio in \code{customPizza}. Alla fine
\code{customPizza} punta all'anello più esterno, e da lì si arriva, di riferimento in riferimento, fino alla pizza base.

\begin{center}
\begin{tikzpicture}
  \memblocco{swblue}{pm}{(0,0)}{2.5cm}{main}{customPizza/{}}
  \memrif{pm-1}
  \foreach \nome/\classe/\x in {p5/Mozzarella/3.4, p4/Fungo/6.4, p3/Mozzarella/9.4, p2/Prosciutto/12.4}{
    \memblocco{swteal}{\nome}{(\x,0)}{2.4cm}{\classe}{pizza/{}}
    \memrif{\nome-1}
  }
  \memblocco{swteal}{p1}{(15.4,0)}{2.6cm}{OlioEPomodoro}{}
  \draw[memfreccia] (pm-1-p) -- (p5-sx);
  \draw[memfreccia] (p5-1-p) -- (p4-sx);
  \draw[memfreccia] (p4-1-p) -- (p3-sx);
  \draw[memfreccia] (p3-1-p) -- (p2-sx);
  \draw[memfreccia] (p2-1-p) -- ++(0.3,0) |- (p1.west);
  % i risultati di getCost() che risalgono la catena
  \foreach \nome/\testo in {p5/{$7{,}50 + 1{,}50 = 9{,}00$}, p4/{$5{,}50 + 2{,}00 = 7{,}50$},
                            p3/{$4{,}00 + 1{,}50 = 5{,}50$}, p2/{$3{,}00 + 1{,}00 = 4{,}00$}, p1/{$3{,}00$}}{
    \node[font=\scriptsize, anchor=north, inner sep=2pt] (r-\nome) at (\nome-giu |- 0,-1.35) {\testo};
  }
  \draw[-{Stealth[length=2mm]}, line width=0.8pt, swred, dashed] (r-p1.south |- 0,-1.95) --
    node[below, font=\scriptsize, text=swred]{\texttt{getCost()}: i risultati risalgono la catena, dalla pizza base
    all'anello più esterno} (r-p5.south |- 0,-1.95);
  \memzona[swblue]{(pm)(pm-fondo)}{Stack}
  \memzona[swteal]{(p5)(p1)(p2-fondo)(r-p5)(r-p1)}{Heap}
\end{tikzpicture}
\end{center}

La chiamata \code{customPizza.getCost()} \textbf{scende} lungo la catena: ogni decoratore, prima di aggiungere il suo
prezzo, chiede il costo alla pizza che contiene. Solo \code{OlioEPomodoro} risponde subito, con 3,00; poi i risultati
\textbf{risalgono}, e ogni anello aggiunge il suo pezzo: 4,00, 5,50, 7,50, 9,00. La descrizione si costruisce allo stesso
modo, ed è per questo che gli ingredienti compaiono nell'ordine in cui sono stati aggiunti.

\info{{Quante classi servirebbero con l'ereditarietà?}}{
  Con tre ingredienti, una sottoclasse per ogni combinazione richiederebbe $2^3 = 8$ classi (\textit{base},
  \textit{Mozzarella}, \textit{Prosciutto}, \textit{Mozzarella e Prosciutto}, \ldots), e non permetterebbe comunque la
  doppia mozzarella; con dieci ingredienti sarebbero $2^{10} = 1024$. Con i decoratori servono tre classi, una per
  ingrediente, e le combinazioni, ripetizioni comprese, si scelgono a tempo di esecuzione.
}

\Needspace{24\baselineskip}
\subsubsection*{Il Decorator nelle librerie di Java}

Il pattern è usato spesso nelle librerie di Java, in particolare nell'input/output: \code{Reader}, \code{BufferedReader(Reader
in)}, \code{InputStream}, \code{InputStreamReader(InputStream in)}, \code{BufferedInputStream(InputStream in)}. Ogni classe
riceve nel costruttore un oggetto del suo stesso tipo (o del tipo da cui legge) e gli aggiunge una capacità:
\begin{itemize}
  \item \code{System.in} è un \code{InputStream}: legge \textbf{byte};
  \item \code{InputStreamReader} lo avvolge e trasforma i byte in \textbf{caratteri}, secondo una codifica;
  \item \code{BufferedReader} avvolge il lettore di caratteri e aggiunge un \textbf{buffer} e il metodo \code{readLine()}, che
        legge una riga intera.
\end{itemize}

\Needspace{30\baselineskip}
\begin{lstlisting}[style=java, name=LeggiRighe.java]
import java.io.BufferedReader;
import java.io.IOException;
import java.io.InputStreamReader;

public class LeggiRighe {
    public static void main(String[] args) throws IOException {
        BufferedReader in = new BufferedReader(new InputStreamReader(System.in));
        String riga;
        while ((riga = in.readLine()) != null) {      // una riga intera alla volta
            System.out.println("Hai scritto: " + riga);
        }
    }
}
\end{lstlisting}
\noindent\begin{minipage}[c]{0.56\textwidth}
\begin{lstlisting}[style=shell, name=Terminale]
$ printf 'ciao\nDesign Patterns\n' | java LeggiRighe.java
Hai scritto: ciao
Hai scritto: Design Patterns
\end{lstlisting}
\end{minipage}\hfill
\begin{minipage}[c]{0.41\textwidth}
\centering
\begin{tikzpicture}
  \memblocco{swteal}{r3}{(0,0)}{3.1cm}{BufferedReader}{in/{}}
  \memrif{r3-1}
  \memblocco{swteal}{r2}{(0,-1.35)}{3.1cm}{InputStreamReader}{in/{}}
  \memrif{r2-1}
  \memblocco{swteal}{r1}{(0,-2.7)}{3.1cm}{InputStream}{}
  \draw[memfreccia] (r3-1-p) -- ++(0.45,0) |- (r2.east);
  \draw[memfreccia] (r2-1-p) -- ++(0.45,0) |- (r1.east);
  \node[font=\scriptsize, anchor=west] at ([xshift=0.7cm]r3.east) {righe};
  \node[font=\scriptsize, anchor=west] at ([xshift=0.7cm]r2.east) {caratteri};
  \node[font=\scriptsize, anchor=west] at ([xshift=0.7cm]r1.east) {byte (\texttt{System.in})};
\end{tikzpicture}
\end{minipage}

\medskip
Nel disegno la catena di oggetti creata dalla prima riga di \code{main}; i nomi dei campi interni delle classi della libreria
sono semplificati.

\Needspace{12\baselineskip}
\subsubsection*{Vantaggi e svantaggi}

\begin{itemize}
  \item \textbf{Vantaggio}: combinando i decoratori si ottiene un numero illimitato di combinazioni di chiamate ai
        costruttori. Non serve dichiarare le classi «prosciutto e funghi», «margherita», «margherita prosciutto e funghi»,
        «margherita doppia mozzarella»\ldots
  \item \textbf{Svantaggio}: costruire un oggetto richiede di istanziare più oggetti e di chiamare più costruttori, e ogni
        chiamata di metodo attraversa tutta la catena. Il pattern è quindi poco efficiente dal punto di vista
        computazionale (e, di conseguenza, anche da quello energetico).
\end{itemize}

% ============================================================
\section{Pattern comportamentali}

I pattern comportamentali riguardano gli \textbf{algoritmi} e l'\textbf{assegnazione delle responsabilità} fra gli oggetti:
come gli oggetti interagiscono e si dividono il lavoro senza legarsi troppo strettamente a un algoritmo particolare. Ne
vediamo tre: Observer e State, a cui le slide dedicano un esempio, e Strategy, che la tabella delle slide evidenzia.

\subsection{Observer}

\dfn{Observer (comportamentale)}{
  Il pattern \term{Observer} definisce una dipendenza 1:N fra oggetti: quando un oggetto (il \textbf{soggetto}) cambia
  stato, tutti gli oggetti che dipendono da lui (gli \textbf{osservatori}) vengono avvisati e aggiornati automaticamente.
}

\begin{itemize}
  \item \textbf{Contesto}: uno o più oggetti sono interessati a osservare i cambiamenti di stato di un soggetto.
  \item \textbf{Problema}: il soggetto può avere osservatori diversi, anche di tipi diversi, e deve essere possibile
        aggiungere nuovi osservatori mentre l'applicazione è in esecuzione.
  \item \textbf{Soluzione}: il soggetto offre delle operazioni con cui un osservatore dichiara il proprio interesse per i
        cambiamenti di stato (\code{attach}) o lo ritira (\code{detach}). Quando lo stato cambia, il soggetto chiama
        \code{notify}, che a sua volta chiama \code{update} su ogni osservatore registrato.
\end{itemize}

\begin{figure}[H]
\centering
\begin{tikzpicture}
  \umlclass[3.0cm]{oas}{(0,0)}{\textit{AbstractSubject}}{}{+ attach(o : Observer)\\+ detach(o : Observer)\\+ notify()}
  \umlclass[2.4cm]{oob}{(6.4,0)}{\textit{Observer}}{}{+ update()}
  \umlclass[3.0cm]{ocs}{(0,-3.0)}{ConcreteSubject}{- state}{+ getState()\\+ setState(s)}
  \umlclass[2.6cm]{oco}{(6.4,-3.0)}{ConcreteObserver}{- observerState}{+ update()}
  \draw[navig] (oas.east) -- node[molt, above]{observers} node[molt, below, pos=0.85]{0..*} (oob.west);
  \draw[gen] (ocs.north) -- (oas.south);
  \draw[gen] (oco.north) -- (oob.south);
  \draw[navig] (oco.west) -- node[molt, above]{subject} node[molt, below, pos=0.85]{1} (ocs.east);
  \node[notauml, anchor=east] (on1) at (-2.3,0.3) {for all o in observers\\\hspace*{1em}o.update()};
  \draw[dashed, line width=0.5pt, draw=black!70] (on1.east) -- ([yshift=-0.25cm]oas.west);
  \node[notauml, anchor=north] (on2) at (6.4,-4.25) {observerState = subject.getState()};
  \draw[dashed, line width=0.5pt, draw=black!70] (on2.north) -- (oco.south);
\end{tikzpicture}
\caption{La struttura dell'Observer. Il soggetto conosce i suoi osservatori solo attraverso l'interfaccia \code{Observer}.}
\end{figure}

\Needspace{16\baselineskip}
Il diagramma di sequenza mostra il comportamento: l'osservatore si registra; più tardi una sorgente esterna cambia lo stato
del soggetto, che notifica tutti gli osservatori; ognuno, ricevuto l'\code{update()}, chiede al soggetto il nuovo stato.

\begin{center}
\begin{tikzpicture}
  \foreach \n/\x/\t in {sq1/0/{\underline{:ConcreteObserver}}, sq2/4.4/{\underline{:ConcreteSubject}}, sq3/8.6/{External Source}}{
    \node[partec, minimum width=2.6cm] (\n) at (\x,0) {\t};
    \draw[vita] (\n.south) -- ++(0,-4.6);
  }
  \draw[sinc] (0.08,-0.9) -- node[msgl]{attach(this)} (4.32,-0.9);
  \attiv{4.4}{-0.9}{-1.2}
  \draw[sinc] (8.52,-1.7) -- node[msgl]{setState(s)} (4.48,-1.7);
  \attiv{8.6}{-1.6}{-1.9}
  \attiv{4.4}{-1.7}{-4.2}
  \draw[sinc] (4.48,-2.2) -- (5.6,-2.2) -- (5.6,-2.5) -- (4.48,-2.5);
  \node[font=\tiny, anchor=west] at (5.65,-2.35) {notify()};
  \draw[sinc] (4.32,-2.9) -- node[msgl]{update()} (0.08,-2.9);
  \attiv{0}{-2.9}{-4.0}
  \draw[sinc] (0.08,-3.35) -- node[msgl]{getState()} (4.32,-3.35);
  \draw[ritorno] (4.32,-3.8) -- node[msgl]{s} (0.08,-3.8);
\end{tikzpicture}
\end{center}

\Needspace{12\baselineskip}
\subsubsection*{L'esempio delle slide}

L'osservatore è un'interfaccia e il soggetto una classe astratta: entrambi sono \textbf{invarianti}, cioè restano uguali
in ogni applicazione, e potrebbero essere un'interfaccia e una classe di libreria. Cambiano solo le classi concrete. I
nomi del codice corrispondono a quelli del diagramma così:

\begin{table}[H]
\centering
\renewcommand{\arraystretch}{1.25}
\begin{tabular}{@{}ll@{}}
\toprule
\textbf{Nel diagramma} & \textbf{Nel codice dell'esempio} \\
\midrule
\code{AbstractSubject}, \code{ConcreteSubject} & \code{Soggetto}, \code{SoggettoConcreto} \\
\code{attach()}, \code{detach()} & \code{registraOsservatore()}, \code{rimuoviOsservatore()} \\
\code{notify()} & \code{update()} di \code{Soggetto} (protetto) \\
\code{Observer}, \code{ConcreteObserver} & \code{Osservatore}, \code{OsservatoreConcreto} \\
\code{update()} dell'osservatore & \code{aggiorna()} \\
\code{observerState} & il campo \code{stato} dell'osservatore \\
\bottomrule
\end{tabular}
\caption{Corrispondenza fra il diagramma e il codice. Il metodo \code{rimuoviOsservatore()} non è nelle slide: l'abbiamo
aggiunto per completare l'interfaccia del soggetto.}
\end{table}

\Needspace{38\baselineskip}
\noindent\begin{minipage}[t]{0.49\textwidth}
\begin{lstlisting}[style=java, name=Osservatore.java]
// osservatore (interfaccia): invariante,
// potrebbe essere un'interfaccia di libreria
public interface Osservatore {
    public void aggiorna();
}
\end{lstlisting}
\begin{lstlisting}[style=java, name=OsservatoreConcreto.java]
// osservatore (concreto): e' legato al soggetto
// di cui vuole conoscere le modifiche
public class OsservatoreConcreto
        implements Osservatore {
    private int stato;           // copia dello stato
    private Soggetto osservato;  // il soggetto

    public OsservatoreConcreto(Soggetto soggetto) {
        this.osservato = soggetto;
    }

    @Override
    public void aggiorna() {
        stato = osservato.getStato();
        System.out.println(
                "Notificato; Nuovo Stato:");
        System.out.println(stato);
    }
}
\end{lstlisting}
\end{minipage}\hfill
\begin{minipage}[t]{0.49\textwidth}
\begin{lstlisting}[style=java, name=Soggetto.java]
import java.util.ArrayList;

// soggetto (astratto): invariante,
// potrebbe essere una classe di libreria
public abstract class Soggetto {
    protected int stato;  // stato del soggetto
    protected ArrayList<Osservatore> osservatori;

    public Soggetto() {
        osservatori = new ArrayList<Osservatore>();
        stato = 0;
    }

    // attach
    public void registraOsservatore(Osservatore o) {
        osservatori.add(o);
    }

    // detach
    public void rimuoviOsservatore(Osservatore o) {
        osservatori.remove(o);
    }

    // notify: avvisa tutti gli osservatori
    protected void update() {
        for (int i = 0; i < osservatori.size(); i++) {
            Osservatore o = osservatori.get(i);
            o.aggiorna();
        }
    }

    public abstract void setStato(int s);
    public abstract int getStato();
}
\end{lstlisting}
\end{minipage}

\Needspace{24\baselineskip}
\noindent\begin{minipage}[t]{0.49\textwidth}
\begin{lstlisting}[style=java, name=Main.java]
public class Main {
    public static void main(String[] args) {
        Soggetto soggetto = new SoggettoConcreto();
        Osservatore o1 =
                new OsservatoreConcreto(soggetto);
        soggetto.registraOsservatore(o1);

        Osservatore o2 =
                new OsservatoreConcreto(soggetto);
        soggetto.registraOsservatore(o2);

        soggetto.setStato(20);
    }
}
\end{lstlisting}
\begin{lstlisting}[style=shell, name=Terminale]
$ javac *.java
$ java Main
Notificato; Nuovo Stato:
20
Notificato; Nuovo Stato:
20
\end{lstlisting}
\end{minipage}\hfill
\begin{minipage}[t]{0.49\textwidth}
\begin{lstlisting}[style=java, name=SoggettoConcreto.java]
// soggetto (concreto)
public class SoggettoConcreto extends Soggetto {
    public SoggettoConcreto() {
        super();
    }

    @Override
    public void setStato(int s) {
        stato = s;
        update();    // la chiamata in piu' rispetto
                     // a una classe normale
    }

    @Override
    public int getStato() {
        return stato;
    }
}
\end{lstlisting}
\end{minipage}

\medskip
Rispetto a una classe qualsiasi, \code{SoggettoConcreto} ha una sola chiamata in più, \code{update()} dentro
\code{setStato}: è quella che fa partire le notifiche. Una sola chiamata a \code{setStato(20)} produce due stampe, una per
osservatore. (Nella slide il testo stampato inizia con una virgoletta «, che in Java non chiude la stringa: qui è
corretta.)

\Needspace{24\baselineskip}
\subsubsection*{Che cosa succede in memoria}

Dopo \code{setStato(20)} la memoria è questa. I riferimenti vanno \textbf{nei due sensi}: il soggetto conosce i suoi
osservatori attraverso la lista (disegnata, come nel Composite, dentro l'oggetto), e ogni osservatore conosce il soggetto
attraverso il campo \code{osservato}, da cui legge il nuovo stato. Ogni osservatore ne tiene una \textbf{copia} nel proprio
campo \code{stato}. Le variabili \code{o1} e \code{o2} di \code{main}, che puntano agli stessi due osservatori, non sono
disegnate.

\begin{center}
\begin{tikzpicture}
  \memblocco{swblue}{om}{(0,0)}{2.6cm}{main}{soggetto/{}}
  \memrif{om-1}
  \memblocco{swteal}{os}{(3.8,0)}{3.6cm}{SoggettoConcreto}{stato/20, {osservatori[0]}/{}, {osservatori[1]}/{}}
  \memrif{os-2}
  \memrif{os-3}
  \memblocco{swteal}{oo1}{(9.2,0.4)}{3.2cm}{OsservatoreConcreto}{stato/20, osservato/{}}
  \memrif{oo1-2}
  \memblocco{swteal}{oo2}{(9.2,-1.95)}{3.2cm}{OsservatoreConcreto}{stato/20, osservato/{}}
  \memrif{oo2-2}
  \draw[memfreccia] (om-1-p) -- ([yshift=-3pt]os.west);
  \draw[memfreccia] (os-2-p) -- ++(0.55,0) |- (oo1.west);
  \draw[memfreccia] (os-3-p) -- ++(0.75,0) |- (oo2.west);
  \coordinate (osopra) at ($(os.north)+(0,0.55)$);
  \coordinate (osotto) at ($(os-giu)+(0,-1.75)$);
  \draw[memfreccia] (oo1-2-p) -- ++(0.55,0) |- (osopra) -- (os.north);
  \draw[memfreccia] (oo2-2-p) -- ++(0.55,0) |- (osotto) -- (os-giu);
  \memzona[swblue]{(om)(om-fondo)}{Stack}
  \memzona[swteal]{(os)(oo1)(oo2-fondo)(osopra)(osotto)}{Heap}
\end{tikzpicture}
\end{center}

\Needspace{14\baselineskip}
\subsubsection*{Un problema: il valore delle azioni}

\noindent\begin{minipage}[c]{0.62\textwidth}
Si vuole gestire un attributo \code{valore} nella classe \code{AzioniAcquistate}, che misuri il valore istantaneo di un
insieme di azioni acquistate di una società. È calcolato come
\[ \text{AzioniAcquistate.valore} = \text{Società.quotazione} \times \text{AzioniAcquistate.quantità}. \]
Quando il sistema aggiorna la quotazione delle azioni di una società, si vuole che gli attributi \code{valore} di tutte le
istanze di \code{AzioniAcquistate} di quella società siano aggiornati \textbf{automaticamente}.
\end{minipage}\hfill
\begin{minipage}[c]{0.35\textwidth}
\centering
\begin{tikzpicture}
  \umlclass[2.4cm]{paz}{(0,0)}{AzioniAcquistate}{+ quantità\\+ valore}{}
  \umlclass[2.0cm]{pso}{(3.6,0)}{Società}{+ quotazione}{}
  \draw[assoc] (paz.east) -- node[molt, above, pos=0.15]{0..*} node[molt, above, pos=0.88]{1} (pso.west);
\end{tikzpicture}
\end{minipage}

\medskip
La soluzione delle slide applica l'Observer: la società è il soggetto, le azioni acquistate sono gli osservatori. Ogni
azione, appena creata, si registra presso la sua società (\code{attach}); \code{setQuotazione} cambia la quotazione e chiama
\code{notify}; l'\code{update} di ogni azione ricalcola \code{valore = societa.getQuotazione() * quantita}.

\begin{figure}[H]
\centering
\begin{tikzpicture}
  \umlclass[3.0cm]{aas}{(0,0)}{\textit{AbstractSocieta}}{}{+ attach(o)\\+ detach(o)\\\# notifica()}
  \umlclass[3.4cm]{aoq}{(6.6,0)}{{\tiny\normalfont «interface»}\\OsservatoreQuotazione}{}{+ update()}
  \umlclass[3.0cm]{aso}{(0,-3.1)}{Societa}{- nome : String\\- quotazione : double}{+ getQuotazione()\\+ setQuotazione(q)}
  \umlclass[3.4cm]{aaz}{(6.6,-3.1)}{AzioneAcquistata}{- quantita : int\\- valore : double}{+ AzioneAcquistata(s, q)\\+ update()}
  \draw[navig] (aas.east) -- node[molt, above]{observers} node[molt, below, pos=0.85]{0..*} (aoq.west);
  \draw[gen] (aso.north) -- (aas.south);
  \draw[realizza] (aaz.north) -- (aoq.south);
  \draw[navig] (aaz.west) -- node[molt, above]{societa} node[molt, below, pos=0.85]{1} (aso.east);
  \node[notauml, anchor=north] (an1) at (0,-4.65) {\{ \ldots\ notifica(); \}};
  \draw[dashed, line width=0.5pt, draw=black!70] (an1.north) -- (aso.south);
  \node[notauml, anchor=north] (an2) at (6.6,-4.65) {\{ valore = societa.getQuotazione() * quantita; \}};
  \draw[dashed, line width=0.5pt, draw=black!70] (an2.north) -- (aaz.south);
  \node[notauml, anchor=west] (an3) at (9.0,-2.3) {\{ societa.attach(this); \}};
  \draw[dashed, line width=0.5pt, draw=black!70] (an3.south west) -- (aaz.north east);
\end{tikzpicture}
\caption{La soluzione con l'Observer. Nelle slide la società concreta si chiama \code{ConcreteSocieta}; qui \code{Societa}.}
\end{figure}

\Needspace{30\baselineskip}
Le slide si fermano al diagramma; ecco il codice.

\noindent\begin{minipage}[t]{0.49\textwidth}
\begin{lstlisting}[style=java, name=OsservatoreQuotazione.java]
public interface OsservatoreQuotazione {
    void update();
}
\end{lstlisting}
\begin{lstlisting}[style=java, name=AbstractSocieta.java]
import java.util.ArrayList;
import java.util.List;

public abstract class AbstractSocieta {
    private final List<OsservatoreQuotazione>
            observers = new ArrayList<>();

    public void attach(OsservatoreQuotazione o) {
        observers.add(o);
    }
    public void detach(OsservatoreQuotazione o) {
        observers.remove(o);
    }

    protected void notifica() {  // il notify()
        for (OsservatoreQuotazione o : observers) {
            o.update();
        }
    }
}
\end{lstlisting}
\end{minipage}\hfill
\begin{minipage}[t]{0.49\textwidth}
\begin{lstlisting}[style=java, name=Societa.java]
public class Societa extends AbstractSocieta {
    private final String nome;
    private double quotazione;

    public Societa(String nome, double quotazione) {
        this.nome = nome;
        this.quotazione = quotazione;
    }

    public String getNome() { return nome; }
    public double getQuotazione() {
        return quotazione;
    }

    public void setQuotazione(double quotazione) {
        this.quotazione = quotazione;
        notifica();    // le azioni si ricalcolano
    }
}
\end{lstlisting}
\end{minipage}

\Needspace{30\baselineskip}
\noindent\begin{minipage}[t]{0.49\textwidth}
\begin{lstlisting}[style=java, name=AzioneAcquistata.java]
public class AzioneAcquistata
        implements OsservatoreQuotazione {
    private final Societa societa;
    private final int quantita;
    private double valore;

    public AzioneAcquistata(Societa societa,
                            int quantita) {
        this.societa = societa;
        this.quantita = quantita;
        societa.attach(this);  // si iscrive
        update();              // valore iniziale
    }

    @Override
    public void update() {
        valore = societa.getQuotazione() * quantita;
    }

    @Override
    public String toString() {
        return quantita + " azioni "
                + societa.getNome() + " = "
                + valore + " euro";
    }
}
\end{lstlisting}
\end{minipage}\hfill
\begin{minipage}[t]{0.49\textwidth}
\begin{lstlisting}[style=java, name=Portafoglio.java]
public class Portafoglio {
    public static void main(String[] args) {
        Societa acme = new Societa("ACME", 10.0);
        AzioneAcquistata mie =
                new AzioneAcquistata(acme, 100);
        AzioneAcquistata tue =
                new AzioneAcquistata(acme, 40);
        System.out.println(mie);
        System.out.println(tue);

        // la borsa aggiorna la quotazione...
        acme.setQuotazione(12.5);
        // ...e i valori sono gia' aggiornati
        System.out.println(mie);
        System.out.println(tue);
    }
}
\end{lstlisting}
\begin{lstlisting}[style=shell, name=Terminale]
$ javac *.java
$ java Portafoglio
100 azioni ACME = 1000.0 euro
40 azioni ACME = 400.0 euro
100 azioni ACME = 1250.0 euro
40 azioni ACME = 500.0 euro
\end{lstlisting}
\end{minipage}

\medskip
Il \code{main} non chiede mai alle azioni di ricalcolarsi: basta cambiare la quotazione, e le due azioni lo vengono a
sapere da sole.

\Needspace{18\baselineskip}
\begin{warningbox}{{Perché notifica() e non notify()}}
  Il diagramma GoF chiama \code{notify()} il metodo che avvisa gli osservatori, ma in Java non si può: \code{notify()}
  esiste già in \code{java.lang.Object}, la superclasse di tutte le classi, serve alla sincronizzazione fra thread ed è
  \code{final}, quindi non si può ridefinire. Con un metodo \code{public void notify() \{ \}} in \code{AbstractSocieta} il
  compilatore rifiuta la classe:
  \begin{lstlisting}[style=shell, name=Terminale]
$ javac AbstractSocieta.java
AbstractSocieta.java:2: error: notify() in AbstractSocieta cannot override notify() in Object
    public void notify() { }
                ^
  overridden method is final
1 error
\end{lstlisting}
  Per questo il codice delle slide usa nomi diversi (\code{update()} nel soggetto, \code{aggiorna()} nell'osservatore), e
  qui usiamo \code{notifica()}.
\end{warningbox}

\Needspace{14\baselineskip}
\subsubsection*{Quando usarlo}

\begin{itemize}
  \item Quando i cambiamenti di un soggetto devono essere «trasmessi» (\textit{broadcast}) a molti oggetti.
  \item Quando oggetti indipendenti dall'applicazione (i soggetti) devono invocare oggetti specifici dell'applicazione per
        notificare un cambiamento del proprio stato interno. Il soggetto non ha bisogno di sapere nulla dei suoi osservatori
        né di come funzionano: per esempio quando un sensore o un servizio fornisce un dato aggiornato e l'interfaccia
        utente deve essere avvisata.
  \item L'Observer è il cuore del pattern \textbf{MVC} (Model-View-Controller), in cui gli oggetti dell'interfaccia utente
        (\textit{View}) possono essere osservatori diretti degli oggetti del modello (\textit{Model}): gli oggetti del
        Model hanno le funzioni di notifica, gli oggetti della GUI le funzioni di aggiornamento.
\end{itemize}

\begin{esempio}{L'Observer nelle interfacce grafiche di Java}
Nelle librerie grafiche di Java (AWT e Swing) ogni pulsante è un soggetto: gli osservatori si chiamano \textit{listener} e
si registrano con i metodi \code{add\ldots Listener}. Quando l'utente fa clic, il pulsante notifica tutti i listener
registrati.
\begin{lstlisting}[style=java, name=Finestra.java]
import javax.swing.JButton;
import javax.swing.JFrame;

public class Finestra {
    public static void main(String[] args) {
        JFrame finestra = new JFrame("Editor");
        JButton salva = new JButton("Salva");
        // attach: il pulsante (soggetto) registra un osservatore
        salva.addActionListener(evento -> System.out.println("Salvataggio..."));
        finestra.add(salva);
        finestra.pack();
        finestra.setVisible(true);
    }
}
\end{lstlisting}
La libreria standard aveva anche una coppia \code{java.util.Observable} / \code{java.util.Observer} pronta all'uso, ma è
deprecata da Java 9: oggi si scrivono interfacce proprie, come negli esempi di questo capitolo, o si usano i
\textit{listener}.
\end{esempio}

\subsection{State}

\dfn{State (comportamentale)}{
  Il pattern \term{State} permette a un oggetto di \textbf{cambiare comportamento a seconda del suo stato interno}. È una
  soluzione generale e flessibile per realizzare le macchine a stati.
}

\noindent\begin{minipage}[c]{0.55\textwidth}
La classe \code{State} contiene implementazioni \textbf{vuote} di tutti i metodi che corrispondono alle transizioni (i
\textit{trigger}, o \textit{handler}). Il polimorfismo permette al \code{Context} di avere un attributo di tipo
\code{State} e di interagire con qualunque stato possibile: \code{request()} esegue \code{state.handle()}, e a seconda
dell'oggetto che c'è in \code{state} viene eseguito il codice di uno stato o dell'altro.
\end{minipage}\hfill
\begin{minipage}[c]{0.42\textwidth}
\centering
\begin{tikzpicture}
  \umlclass[1.9cm]{sct}{(0,0)}{Context}{}{+ request()}
  \umlclass[1.9cm]{sst}{(3.6,0)}{\textit{State}}{}{+ handle()}
  \umlclass[2.4cm]{ssa}{(2.1,-2.2)}{ConcreteStateA}{}{+ handle()}
  \umlclass[2.4cm]{ssb}{(5.1,-2.2)}{ConcreteStateB}{}{+ handle()}
  \draw[aggr] (sct.east) -- node[molt, above]{state} (sst.west);
  \draw[gen] (ssa.north) -- ++(0,0.45) -| (sst.south);
  \draw[gen] (ssb.north) -- ++(0,0.45) -| (sst.south);
  \node[notauml, anchor=north] (scn) at (-0.3,-1.4) {state.handle()};
  \draw[dashed, line width=0.5pt, draw=black!70] (scn.north) -- (sct.south);
\end{tikzpicture}
\end{minipage}

\medskip
Il pattern l'abbiamo già incontrato nel Capitolo~9, dove traduceva in codice il diagramma di stato del
\textbf{distributore automatico}: là si trovano il diagramma di stato, la classe contenitore \code{Distributore}, la classe
generale \code{State} e lo stato \code{Attesa}. Il diagramma di stato si traduce così:
\begin{itemize}
  \item una \textbf{classe generale} che dichiara i metodi corrispondenti a tutte le azioni, ma li implementa senza fare
        nulla;
  \item una serie di \textbf{classi derivate}, una per ogni stato del diagramma, che implementano (con un
        \textit{override}) solo i metodi corrispondenti agli eventi che possono scattare in quello stato.
\end{itemize}
Qui completiamo l'esempio. Il diagramma delle classi delle slide mostra quali eventi gestisce ciascuno stato.

\begin{figure}[H]
\centering
\begin{tikzpicture}
  \node[classe1=1.6cm] (zcl) at (0,3.9) {Cliente};
  \umlclass[5.0cm]{zdi}{(0,0)}{Distributore}
    {- stato : State\\- credito : int\\\underline{\textasciitilde\ attesa : State}\\\underline{\textasciitilde\ attesaSceltaProdotto : State}\\\underline{\textasciitilde\ erogazioneProdotto : State}}
    {+ inserisciMoneta()\\+ selezionaProdotto()\\+ richiestaRestituzioneCredito()\\+ prelevaProdotto()\\+ timeoutErogazione()\\+ spegniDistributore()\\\textasciitilde\ setState(s : State)\\\textasciitilde\ getCredito() : int\\\textasciitilde\ setCredito(c : int)}
  \umlclass[4.2cm]{zst}{(7.6,0)}{State}{}{+ inserisciMoneta(d)\\+ sceltaProdotto(d)\\+ richiestaRestituzione(d)\\+ prelevaProdotto(d)\\+ timeOutErogazione(d)\\+ spegniDistributore(d)}
  \draw[dep] (zcl.south) -- node[stereo, right]{«create»} (zdi.north);
  \draw[compos] (zdi.east) -- node[molt, above]{stato} node[molt, below, pos=0.88]{1} (zst.west);
  \umlclass[4.0cm]{zs1}{(0.4,-5.4)}{Attesa}{}{+ inserisciMoneta(d)\\+ sceltaProdotto(d)\\+ richiestaRestituzione(d)\\+ prelevaProdotto(d)\\+ spegniDistributore(d)}
  \umlclass[4.0cm]{zs2}{(5.0,-5.4)}{AttesaSceltaProdotto}{}{+ inserisciMoneta(d)\\+ sceltaProdotto(d)\\+ richiestaRestituzione(d)\\+ prelevaProdotto(d)\\+ spegniDistributore(d)}
  \umlclass[4.0cm]{zs3}{(9.6,-5.4)}{ErogazioneProdotto}{}{+ inserisciMoneta(d)\\+ sceltaProdotto(d)\\+ richiestaRestituzione(d)\\+ prelevaProdotto(d)\\+ timeOutErogazione(d)\\+ spegniDistributore(d)}
  \foreach \s in {zs1, zs2, zs3}{ \draw[gen] (\s.north) |- (7.6,-3.0) -- (zst.south); }
\end{tikzpicture}
\caption{Il distributore con il pattern State, dalle slide (\code{d} è il \code{Distributore} passato a ogni metodo). Il
simbolo \textasciitilde\ indica la visibilità di package; i tre stati sono attributi statici, quindi sottolineati.}
\end{figure}

\Needspace{30\baselineskip}
Le slide mostrano il codice di \code{Distributore}, di \code{State} e di \code{Attesa} (quello del Capitolo~9) e dicono che
la soluzione completa è su Teams. Qui ricostruiamo gli altri due stati dal diagramma di stato del Capitolo~9: una moneta in
più in \textbf{Attesa scelta prodotto} incrementa il credito; la scelta del prodotto avvia l'\textbf{erogazione};
l'erogazione avvenuta (il prelievo) scala 2 euro e porta in Attesa scelta prodotto se prima il credito era di almeno 4 euro,
altrimenti in Attesa; il \textit{time out} torna alla scelta senza scalare nulla. I messaggi stampati sono nostri.

\begin{lstlisting}[style=java, name=AttesaSceltaProdotto.java]
public class AttesaSceltaProdotto extends State {
    @Override
    public void inserisciMoneta(Distributore distributore) {      // autotransizione
        distributore.setCredito(distributore.getCredito() + 1);
        System.out.println("Credito attuale: " + distributore.getCredito());
    }
    @Override
    public void richiestaRestituzione(Distributore distributore) {
        System.out.println("Ritira il tuo credito: " + distributore.getCredito());
        distributore.setCredito(0);
        distributore.setState(Distributore.attesa);
    }
    @Override
    public void sceltaProdotto(Distributore distributore) {
        System.out.println("Erogazione in corso...");
        distributore.setState(Distributore.erogazioneProdotto);
    }
    @Override
    public void spegniDistributore(Distributore distributore) {
        System.out.println("Ritira il tuo credito: " + distributore.getCredito());
        distributore.setCredito(0);
        System.out.println("Spengo il distributore");
        System.exit(0);
    }
    @Override
    public void prelevaProdotto(Distributore distributore) {
        System.out.println("Il prodotto non e' pronto");
    }
}
\end{lstlisting}

\Needspace{40\baselineskip}
\begin{lstlisting}[style=java, name=ErogazioneProdotto.java]
public class ErogazioneProdotto extends State {
    @Override
    public void prelevaProdotto(Distributore distributore) {     // erogazione avvenuta
        int credito = distributore.getCredito();
        distributore.setCredito(credito - 2);
        System.out.println("Ecco il prodotto. Credito residuo: " + distributore.getCredito());
        if (credito >= 4) {
            distributore.setState(Distributore.attesaSceltaProdotto);
        } else {
            distributore.setState(Distributore.attesa);
        }
    }
    @Override
    public void timeOutErogazione(Distributore distributore) {   // erogazione fallita
        System.out.println("Erogazione fallita: scegli di nuovo");
        distributore.setState(Distributore.attesaSceltaProdotto);
    }
    @Override
    public void inserisciMoneta(Distributore distributore) {
        System.out.println("Attendi: erogazione in corso");
    }
    @Override
    public void sceltaProdotto(Distributore distributore) {
        System.out.println("Attendi: erogazione in corso");
    }
    @Override
    public void richiestaRestituzione(Distributore distributore) {
        System.out.println("Attendi: erogazione in corso");
    }
    @Override
    public void spegniDistributore(Distributore distributore) {
        System.out.println("Ritira il tuo credito: " + distributore.getCredito());
        distributore.setCredito(0);
        System.out.println("Spengo il distributore");
        System.exit(0);
    }
}
\end{lstlisting}

\Needspace{30\baselineskip}
Il \code{main} di prova delle slide simula un cliente. Nella stessa sequenza di comandi la stessa chiamata,
\code{dis.selezionaProdotto()}, produce effetti diversi a seconda dello stato:

\begin{lstlisting}[style=java, name=Cliente.java]
public class Cliente {
    public static void main(String[] args) {
        Distributore dis = new Distributore();

        // scelgo un prodotto senza aver inserito monete
        System.out.print("Scelgo un prodotto. "); dis.selezionaProdotto();
        System.out.print("Inserisco un euro. "); dis.inserisciMoneta();
        System.out.print("Scelgo un prodotto. "); dis.selezionaProdotto();
        System.out.print("Inserisco un euro. "); dis.inserisciMoneta();
        System.out.print("Scelgo un prodotto. "); dis.selezionaProdotto();
        System.out.print("Prelevo il prodotto. "); dis.prelevaProdotto();
        System.out.print("Inserisco un euro. "); dis.inserisciMoneta();
        System.out.print("Inserisco un euro. "); dis.inserisciMoneta();
        System.out.print("Inserisco un euro. "); dis.inserisciMoneta();
        System.out.print("Inserisco un euro. "); dis.inserisciMoneta();
        System.out.print("Scelgo un prodotto. "); dis.selezionaProdotto();
        System.out.print("Scelgo un prodotto. "); dis.selezionaProdotto();
        System.out.print("Richiedo la restituzione del credito. ");
        dis.richiestaRestituzioneCredito();
        System.out.print("Inserisco un euro. "); dis.inserisciMoneta();
        System.out.print("Prelevo il prodotto. "); dis.prelevaProdotto();
        System.out.print("Spengo il distributore. "); dis.spegniDistributore();
        System.out.print("Inserisco un euro. "); dis.inserisciMoneta();
    }
}
\end{lstlisting}

\Needspace{22\baselineskip}
\begin{lstlisting}[style=shell, name=Terminale]
$ javac *.java
$ java Cliente
Scelgo un prodotto. Non hai credito sufficiente
Inserisco un euro. Credito attuale: 1
Scelgo un prodotto. Non hai credito sufficiente
Inserisco un euro. Credito attuale: 2
Scelgo un prodotto. Erogazione in corso...
Prelevo il prodotto. Ecco il prodotto. Credito residuo: 0
Inserisco un euro. Credito attuale: 1
Inserisco un euro. Credito attuale: 2
Inserisco un euro. Credito attuale: 3
Inserisco un euro. Credito attuale: 4
Scelgo un prodotto. Erogazione in corso...
Scelgo un prodotto. Attendi: erogazione in corso
Richiedo la restituzione del credito. Attendi: erogazione in corso
Inserisco un euro. Attendi: erogazione in corso
Prelevo il prodotto. Ecco il prodotto. Credito residuo: 2
Spengo il distributore. Ritira il tuo credito: 2
Spengo il distributore
\end{lstlisting}

La prima scelta trova il distributore in Attesa (\textit{non hai credito sufficiente}), la terza in Attesa scelta prodotto
(parte l'erogazione), la seconda delle due scelte consecutive in Erogazione prodotto (\textit{attendi}). L'ultima riga del
\code{main}, «Inserisco un euro», non viene mai eseguita: \code{spegniDistributore()} termina il programma con
\code{System.exit(0)}.

\Needspace{24\baselineskip}
\subsubsection*{Che cosa succede in memoria}

Il disegno mostra la memoria prima dello spegnimento. Gli stati sono tre oggetti creati \textbf{una volta sola}, nei campi
statici di \code{Distributore}; il campo \code{stato} del distributore punta a uno di loro. Una \textbf{transizione} non crea
nulla: \code{setState(\ldots)} sposta il riferimento. L'ultimo prelievo, con 4 euro di credito, l'ha spostato da
ErogazioneProdotto ad AttesaSceltaProdotto (la freccia grigia è quella di prima).

\begin{center}
\begin{tikzpicture}
  \memblocco{swblue}{zm}{(0,0)}{2.2cm}{main}{dis/{}}
  \memrif{zm-1}
  \memblocco{swteal}{zd}{(3.3,0)}{2.7cm}{Distributore}{stato/{}, credito/2}
  \memrif{zd-1}
  \memblocco{swteal}{za}{(8.3,0.75)}{3.5cm}{Attesa}{}
  \memblocco{swteal}{zb}{(8.3,-0.55)}{3.5cm}{AttesaSceltaProdotto}{}
  \memblocco{swteal}{zc}{(8.3,-1.85)}{3.5cm}{ErogazioneProdotto}{}
  \memblocco{ph3}{zs}{(4.6,-3.4)}{4.4cm}{Distributore}{attesa/{}, attesaSceltaProdotto/{}, erogazioneProdotto/{}}
  \memrif{zs-1}
  \memrif{zs-2}
  \memrif{zs-3}
  \draw[memfreccia] (zm-1-p) -- ([yshift=-3pt]zd.west);
  \draw[memfreccia, draw=black!30, dashed] (zd-1-p) -- ++(0.75,0) |- (zc.west);
  \draw[memfreccia, draw=swred, line width=1.1pt] (zd-1-p) -- ++(0.75,0) |- (zb.west);
  \node[font=\tiny, text=swred, anchor=south east] at ([xshift=-2pt]zb.north west) {stato corrente};
  \draw[memfreccia] (zs-1-p) -- (12.25,0 |- zs-1-p) |- (za.east);
  \draw[memfreccia] (zs-2-p) -- (12.55,0 |- zs-2-p) |- (zb.east);
  \draw[memfreccia] (zs-3-p) -- (12.85,0 |- zs-3-p) |- (zc.east);
  \memzona[swblue]{(zm)(zm-fondo)}{Stack}
  \memzona[swteal]{(zd)(za)(zc-fondo)}{Heap}
  \memzona[ph3]{(zs)(zs-fondo)}{Campi statici}
\end{tikzpicture}
\end{center}

\info{{Perché gli stati possono essere condivisi}}{
  Gli oggetti stato \textbf{non hanno campi}: tutti i dati (il credito) stanno nel contesto, che ogni metodo riceve come
  parametro. Per questo bastano tre oggetti per tutti i distributori del programma, creati una volta sola nei campi
  statici: due distributori diversi possono trovarsi nello stesso stato condividendo lo stesso oggetto, ciascuno con il
  proprio credito. Se uno stato dovesse ricordare qualcosa di suo, servirebbe un oggetto stato per ogni distributore.
}

\subsection{Strategy}

\dfn{Strategy (comportamentale)}{
  Il pattern \term{Strategy} definisce una famiglia di algoritmi, incapsula ciascuno in una classe e li rende
  \textbf{intercambiabili}: l'algoritmo può variare indipendentemente da chi lo usa.
}

Le slide lo evidenziano nella tabella dei pattern comportamentali senza dedicargli un esempio. Eccone uno breve, anche
perché la sua struttura è identica a quella dello State, e conviene vedere in che cosa i due differiscono. Un ordine deve
calcolare il costo della spedizione, e i modi di spedire sono diversi: ciascuno diventa una classe che implementa la stessa
interfaccia.

\begin{center}
\begin{tikzpicture}
  \umlclass[4.4cm]{yor}{(-0.6,0)}{Ordine}{- spedizione : StrategiaSpedizione}{+ setSpedizione(s)\\+ totale() : double}
  \umlclass[3.4cm]{yst}{(5.6,0)}{{\tiny\normalfont «interface»}\\StrategiaSpedizione}{}{+ costo(pesoKg) : double}
  \umlclass[2.2cm]{ys1}{(2.9,-2.6)}{Standard}{}{+ costo(pesoKg)}
  \umlclass[2.2cm]{ys2}{(5.6,-2.6)}{Express}{}{+ costo(pesoKg)}
  \umlclass[2.4cm]{ys3}{(8.4,-2.6)}{RitiroInNegozio}{}{+ costo(pesoKg)}
  \draw[aggr] (yor.east) -- node[molt, above]{strategia} (yst.west);
  \foreach \s in {ys1, ys2, ys3}{ \draw[realizza] (\s.north) |- (5.6,-1.45) -- (yst.south); }
\end{tikzpicture}
\end{center}

\Needspace{26\baselineskip}
\noindent\begin{minipage}[t]{0.49\textwidth}
\begin{lstlisting}[style=java, name=StrategiaSpedizione.java]
// Strategy: la famiglia di algoritmi,
// dietro un'unica interfaccia
public interface StrategiaSpedizione {
    double costo(double pesoKg);
}
\end{lstlisting}
\begin{lstlisting}[style=java, name=Standard.java]
public class Standard implements StrategiaSpedizione {
    public double costo(double pesoKg) {
        return 5.0 + 0.5 * pesoKg;
    }
}
\end{lstlisting}
\begin{lstlisting}[style=java, name=Express.java]
public class Express implements StrategiaSpedizione {
    public double costo(double pesoKg) {
        return 12.0 + 1.0 * pesoKg;
    }
}
\end{lstlisting}
\begin{lstlisting}[style=java, name=RitiroInNegozio.java]
public class RitiroInNegozio
        implements StrategiaSpedizione {
    public double costo(double pesoKg) {
        return 0.0;
    }
}
\end{lstlisting}
\end{minipage}\hfill
\begin{minipage}[t]{0.49\textwidth}
\begin{lstlisting}[style=java, name=Ordine.java]
// Context: usa una strategia senza sapere quale
public class Ordine {
    private final double merce;
    private final double pesoKg;
    private StrategiaSpedizione spedizione =
            new Standard();

    public Ordine(double merce, double pesoKg) {
        this.merce = merce;
        this.pesoKg = pesoKg;
    }

    // la sceglie il client
    public void setSpedizione(StrategiaSpedizione s) {
        this.spedizione = s;
    }

    public double totale() {
        return merce + spedizione.costo(pesoKg);
    }
}
\end{lstlisting}
\end{minipage}

\Needspace{16\baselineskip}
\noindent\begin{minipage}[t]{0.62\textwidth}
\begin{lstlisting}[style=java, name=Negozio.java]
public class Negozio {
    public static void main(String[] args) {
        // 40 euro di merce, 6 kg
        Ordine ordine = new Ordine(40.0, 6.0);
        System.out.println("Standard:   " + ordine.totale());

        // si cambia algoritmo
        ordine.setSpedizione(new Express());
        System.out.println("Express:    " + ordine.totale());

        ordine.setSpedizione(new RitiroInNegozio());
        System.out.println("In negozio: " + ordine.totale());
    }
}
\end{lstlisting}
\end{minipage}\hfill
\begin{minipage}[t]{0.35\textwidth}
\begin{lstlisting}[style=shell, name=Terminale]
$ javac *.java
$ java Negozio
Standard:   48.0
Express:    58.0
In negozio: 40.0
\end{lstlisting}
\end{minipage}

\begin{table}[H]
\centering
\renewcommand{\arraystretch}{1.3}
\begin{tabularx}{\textwidth}{>{\bfseries}P{3.4cm} L L}
\toprule
 & \textbf{State} & \textbf{Strategy} \\
\midrule
Chi cambia l'oggetto & Gli stati stessi, con \code{setState}, quando scatta una transizione & Il client, con
  \code{setSpedizione}, quando vuole \\
Gli oggetti si conoscono? & Sì: ogni stato sa in quali altri stati si può passare & No: ogni strategia ignora le altre \\
Che cosa rappresenta & Una fase della vita dell'oggetto & Un modo diverso di fare la stessa cosa \\
\bottomrule
\end{tabularx}
\caption{Stessa struttura, intenzioni diverse.}
\end{table}

Anche questo pattern è nella libreria di Java: un \code{Comparator} è una strategia di confronto, e la stessa chiamata
\code{lista.sort(c)} ordina la lista per prezzo, per nome o per data a seconda del \code{Comparator c} che le si passa.

% ============================================================
\section{Pattern moderni: la Dependency Injection}

Non tutti i pattern usati oggi fanno parte del catalogo originale della GoF. Uno dei più diffusi, e già incontrato in questa
dispensa, è la Dependency Injection.

\dfn{Dependency Injection}{
  Si ha \term{Dependency Injection} (iniezione delle dipendenze) quando le risorse da cui un componente dipende gli vengono
  \textbf{passate (iniettate) dall'esterno}, invece di essere istanziate dal componente stesso durante la sua esecuzione. Il
  componente dichiara, per esempio nel costruttore, il tipo (eventualmente generico) e la quantità di tutte le risorse da
  cui dipende; i riferimenti alle risorse gli vengono poi passati, per esempio attraverso il costruttore.
}

L'avevamo vista all'opera nel Capitolo~4, dove la \code{Cassa} riceveva nel costruttore il \code{GatewayPagamenti} (così da
poterla provare con un gateway finto), e nel Capitolo~10, dove il \code{Traveler} riceveva dall'esterno il mezzo di
trasporto.

\Needspace{20\baselineskip}
\subsubsection*{Un esempio semplice}

\noindent\begin{minipage}[t]{0.49\textwidth}
\begin{lstlisting}[style=java, name=MyClass.java]
import java.util.logging.Logger;

public class MyClass {
    private Logger logger;

    public MyClass(Logger logger) {
        this.logger = logger;
        // scrive un messaggio informativo nel log
        logger.info("This is a log message.");
    }
}
\end{lstlisting}
\end{minipage}\hfill
\begin{minipage}[t]{0.49\textwidth}
\begin{lstlisting}[style=java, name=Avvio.java]
import java.util.logging.Logger;

public class Avvio {
    public static void main(String[] args) {
        // prima si crea la dipendenza...
        Logger log = Logger.getLogger("ordini");
        // ...poi la si inietta
        MyClass oggetto = new MyClass(log);
    }
}
\end{lstlisting}
\begin{lstlisting}[style=shell, name=Terminale]
$ javac *.java
$ java Avvio
Oct 25, 2026 10:15:39 PM MyClass <init>
INFO: This is a log message.
\end{lstlisting}
\end{minipage}

\medskip
\code{MyClass} può scrivere in un log, e il log su cui scrive è l'oggetto di tipo \code{Logger} che riceve come parametro: non
sa (e non deve sapere) dove finiscono i messaggi. Qui usiamo il \code{Logger} della libreria standard,
\code{java.util.logging}, che scrive su standard error la data, la classe e il metodo che hanno prodotto il messaggio
(\code{<init>} è il costruttore). Nella slide alla chiamata di \code{logger.info} manca il punto e virgola. L'esempio viene
dal tutorial di Vogella sulla Dependency Injection (\url{https://www.vogella.com/tutorials/DependencyInjection/article.html}).

\Needspace{12\baselineskip}
\subsubsection*{Le conseguenze}

\begin{itemize}
  \item Rende la classe \textbf{più generica}, perché ne separa il comportamento dall'oggetto specifico che le viene passato.
        Resta però una dipendenza fra l'oggetto e una classe (o un'interfaccia): il \textbf{tipo} della risorsa.
  \item Non è sempre ottimale dal punto di vista dell'\textbf{accoppiamento}: chi crea l'oggetto deve conoscere e creare anche
        tutte le sue dipendenze.
  \item Si chiama anche \term{inversione del controllo} (\textit{inversion of control}), perché richiede di creare
        \textbf{prima} l'oggetto da cui si dipende e \textbf{poi} quello che lo userà: la direzione è opposta, per esempio, a
        quella del modello BCE (Boundary-Control-Entity), in cui l'interfaccia (Boundary) crea il controllo che le serve e il
        controllo crea e usa le entità.
\end{itemize}

\Needspace{16\baselineskip}
\subsubsection*{Dependency Injection e test}

La Dependency Injection è fondamentale per il \textbf{test in isolamento} e per il \textbf{test di integrazione}: permette di
provare un oggetto passandogli un oggetto \textbf{fittizio} su cui agire, al posto dell'oggetto concreto, che potrebbe non
essere disponibile o non essere ancora stato testato. Il caso tipico delle slide: vogliamo testare una classe che chiama un
servizio, ma il servizio non è ancora disponibile.
\begin{itemize}
  \item Se si segue il pattern, la classe sotto test riceve attraverso il costruttore un riferimento all'oggetto servizio:
        nel test le si passa un servizio fittizio, e la classe si può eseguire anche senza quello originale.
  \item Se non lo si segue, il riferimento al servizio è scritto dentro la classe (oppure il servizio viene istanziato al
        suo interno) e al momento del test non si può cambiare.
\end{itemize}

\Needspace{30\baselineskip}
\noindent\begin{minipage}[t]{0.49\textwidth}
\begin{lstlisting}[style=java, name=ServizioMeteo.java]
public interface ServizioMeteo {
    double temperatura(String citta);   // gradi C
}
\end{lstlisting}
\begin{lstlisting}[style=java, name=AllarmeCaldo.java]
public class AllarmeCaldo {
    private final ServizioMeteo meteo;   // la risorsa

    public AllarmeCaldo(ServizioMeteo meteo) {
        this.meteo = meteo;              // iniettata
    }

    public String controlla(String citta) {
        double t = meteo.temperatura(citta);
        if (t >= 35) {
            return "ALLARME: a " + citta
                    + " ci sono " + t + " gradi";
        }
        return "Nessun allarme a " + citta;
    }
}
\end{lstlisting}
\end{minipage}\hfill
\begin{minipage}[t]{0.49\textwidth}
\begin{lstlisting}[style=java, name=MeteoFinto.java]
// Oggetto fittizio per i test: restituisce
// sempre la temperatura scelta dal test
public class MeteoFinto implements ServizioMeteo {
    private final double temperatura;

    public MeteoFinto(double temperatura) {
        this.temperatura = temperatura;
    }

    @Override
    public double temperatura(String citta) {
        return temperatura;
    }
}
\end{lstlisting}
\end{minipage}
\begin{lstlisting}[style=java, name=App.java]
public class App {
    public static void main(String[] args) {
        // prima la dipendenza...
        ServizioMeteo meteo = new MeteoHttp("https://api.meteo.example");
        // ...poi chi la usa
        AllarmeCaldo allarme = new AllarmeCaldo(meteo);
        System.out.println(allarme.controlla("Napoli"));
    }
}
\end{lstlisting}

\Needspace{30\baselineskip}
In produzione \code{AllarmeCaldo} riceve \code{MeteoHttp}, che interroga un servizio REST con l'\code{HttpClient} visto nel
Capitolo~14; nel test riceve \code{MeteoFinto}, che risponde subito e sempre allo stesso modo. Il test non ha bisogno né
della rete né del servizio, e può scegliere la temperatura che gli serve per provare entrambi i casi.

\begin{lstlisting}[style=java, name=MeteoHttp.java]
import java.net.URI;
import java.net.http.HttpClient;
import java.net.http.HttpRequest;
import java.net.http.HttpResponse;

// Implementazione vera: interroga un servizio REST esterno
public class MeteoHttp implements ServizioMeteo {
    private final HttpClient client = HttpClient.newHttpClient();
    private final String baseUrl;

    public MeteoHttp(String baseUrl) {
        this.baseUrl = baseUrl;
    }

    @Override
    public double temperatura(String citta) {
        try {
            URI uri = URI.create(baseUrl + "/temperatura?citta=" + citta);
            HttpRequest richiesta = HttpRequest.newBuilder(uri).build();
            String corpo = client.send(richiesta, HttpResponse.BodyHandlers.ofString()).body();
            return Double.parseDouble(corpo);
        } catch (Exception e) {
            throw new IllegalStateException("Servizio meteo non raggiungibile", e);
        }
    }
}
\end{lstlisting}

\Needspace{34\baselineskip}
\begin{lstlisting}[style=java, name=AllarmeCaldoTest.java]
import static org.junit.jupiter.api.Assertions.assertEquals;
import org.junit.jupiter.api.Test;

class AllarmeCaldoTest {
    @Test
    void segnalaIlCaldoDa35GradiInSu() {
        AllarmeCaldo allarme = new AllarmeCaldo(new MeteoFinto(38.0));
        assertEquals("ALLARME: a Napoli ci sono 38.0 gradi", allarme.controlla("Napoli"));
    }

    @Test
    void taceSottoI35Gradi() {
        AllarmeCaldo allarme = new AllarmeCaldo(new MeteoFinto(22.5));
        assertEquals("Nessun allarme a Napoli", allarme.controlla("Napoli"));
    }
}
\end{lstlisting}
\begin{lstlisting}[style=shell, name=Terminale, deletestring={[b]'}]
$ java -jar junit-platform-console-standalone.jar execute --class-path out \
       --select-class AllarmeCaldoTest --details=tree --details-theme=ascii
+-- JUnit Jupiter [OK]
| '-- AllarmeCaldoTest [OK]
|   +-- taceSottoI35Gradi() [OK]
|   '-- segnalaIlCaldoDa35GradiInSu() [OK]
...
[         2 tests successful      ]
[         0 tests failed          ]
\end{lstlisting}

\medskip
(L'uscita di JUnit è accorciata: le righe omesse riassumono i contenitori di test trovati ed eseguiti.)

\Needspace{24\baselineskip}
\subsubsection*{Che cosa succede in memoria}

La classe \code{AllarmeCaldo} è la stessa nei due casi; cambia solo l'oggetto a cui punta il suo campo \code{meteo}, e lo
decide chi crea l'allarme.

\noindent\begin{minipage}[t]{0.49\textwidth}
\centering
{\small\textbf{In produzione} (\code{App.main})\par}
\medskip
\begin{tikzpicture}
  \memblocco{swblue}{dm}{(0,0)}{2.0cm}{main}{allarme/{}}
  \memrif{dm-1}
  \memblocco{swteal}{da}{(2.9,0)}{2.4cm}{AllarmeCaldo}{meteo/{}}
  \memrif{da-1}
  \memblocco{swteal}{dh}{(2.9,-1.6)}{5.6cm}{MeteoHttp}{client/{\ldots}, baseUrl/{"https://api.meteo.example"}}
  \draw[memfreccia] (dm-1-p) -- ([yshift=-3pt]da.west);
  \draw[memfreccia] (da-1-p) -- ++(0.45,0) coordinate (dgomito) -- (dgomito |- dh.north);
  \memzona[swblue]{(dm)(dm-fondo)}{Stack}
  \memzona[swteal]{(da)(dh-fondo)}{Heap}
\end{tikzpicture}
\end{minipage}\hfill
\begin{minipage}[t]{0.49\textwidth}
\centering
{\small\textbf{Nel test} (\code{segnalaIlCaldoDa35GradiInSu})\par}
\medskip
\begin{tikzpicture}
  \memblocco{swblue}{tm}{(0,0)}{2.0cm}{test}{allarme/{}}
  \memrif{tm-1}
  \memblocco{swteal}{ta}{(2.9,0)}{2.4cm}{AllarmeCaldo}{meteo/{}}
  \memrif{ta-1}
  \memblocco{swteal}{tf}{(2.9,-1.6)}{2.8cm}{MeteoFinto}{temperatura/38.0}
  \draw[memfreccia] (tm-1-p) -- ([yshift=-3pt]ta.west);
  \draw[memfreccia] (ta-1-p) -- ++(0.45,0) coordinate (tgomito) -- (tgomito |- tf.north);
  \memzona[swblue]{(tm)(tm-fondo)}{Stack}
  \memzona[swteal]{(ta)(tf-fondo)}{Heap}
\end{tikzpicture}
\end{minipage}

\Needspace{16\baselineskip}
\begin{esempio}{Il controesempio: la dipendenza scritta dentro la classe}
\begin{lstlisting}[style=java, name=AllarmeCaldoRigido.java]
public class AllarmeCaldoRigido {
    // la dipendenza e' creata qui dentro: nessun test potra' sostituirla
    private final ServizioMeteo meteo = new MeteoHttp("https://api.meteo.example");

    public String controlla(String citta) {
        double t = meteo.temperatura(citta);      // chiama sempre il servizio vero
        return t >= 35 ? "ALLARME: a " + citta + " ci sono " + t + " gradi"
                       : "Nessun allarme a " + citta;
    }
}
\end{lstlisting}
Per provare \code{AllarmeCaldoRigido} serve il servizio vero, raggiungibile e acceso; e non c'è modo di fargli restituire
38 gradi a comando per verificare l'allarme.
\end{esempio}

\info{{Chi crea gli oggetti? I framework}}{
  In un'applicazione grande, creare a mano tutte le dipendenze nell'ordine giusto diventa faticoso. Framework come
  \textbf{Spring} o \textbf{Jakarta EE} lo fanno da soli: si dichiara nel costruttore di che cosa ha bisogno una classe (in
  Jakarta EE con l'annotazione \code{@Inject}; in Spring basta il costruttore, o \code{@Autowired}), e il framework crea gli
  oggetti e li collega. Il codice delle classi resta quello visto qui.
}

% ============================================================
\section{Conclusioni}

\begin{table}[H]
\centering
\renewcommand{\arraystretch}{1.35}
\begin{tabularx}{\textwidth}{>{\bfseries}P{3.2cm} L L}
\toprule
Famiglia & Su che cosa si concentra & Quali problemi risolve \\
\midrule
Strutturali & Come gli oggetti si compongono per formare strutture più grandi & Realizzare funzionalità nuove a partire
  da quelle esistenti; dare flessibilità ed estensibilità \\
Comportamentali & Gli algoritmi e l'assegnazione delle responsabilità agli oggetti & L'accoppiamento troppo stretto con un
  algoritmo particolare \\
Creazionali & La creazione di oggetti complessi & Nascondere come gli oggetti complessi vengono creati e assemblati \\
\bottomrule
\end{tabularx}
\caption{Le tre famiglie di pattern in sintesi.}
\end{table}

\Needspace{10\baselineskip}
Alcune osservazioni finali:
\begin{itemize}
  \item I pattern permettono di progettare a un \textbf{livello più astratto}: trattano molte interazioni fra classi e oggetti
        come un'unità; spesso sono utili dopo un primo progetto, come obiettivi di ristrutturazione delle classi
        (\textit{refactoring}).
  \item Una progettazione \textbf{orientata alla variazione}: si considera quali aspetti del progetto possono variare, si
        individuano i pattern applicabili, li si varia per valutarne i compromessi, e si ripete.
  \item I design pattern forniscono soluzioni a problemi comuni, portano a modelli e codice estensibili, si possono usare così
        come sono oppure come esempi di \textbf{ereditarietà di interfaccia} e di \textbf{delega}, e applicano gli stessi
        principi alla struttura e al comportamento.
\end{itemize}

\warning{I design pattern non risolvono tutti i problemi}{
  Le slide chiudono con una provocazione: «i design pattern risolvono tutti i vostri problemi di ingegneria del software».
  Non è così. Si potrebbe continuare a lungo a presentare pattern, ma alla fine bisogna \textbf{prima definire il problema}
  e poi cercare una soluzione; e poiché i pattern sono soluzioni a problemi, \textbf{non si guardano i pattern finché il
  problema non è stato definito}. Un pattern applicato dove non serve aggiunge solo classi e complessità.
}

\Needspace{14\baselineskip}
\subsubsection*{Per approfondire}

\begin{itemize}
  \item C. Alexander, \textit{The Timeless Way of Building}, Oxford, 1979; \textit{A Pattern Language}, Oxford, 1977.
  \item E. Gamma, R. Helm, R. Johnson, J. Vlissides, \textit{Design Patterns}, Addison-Wesley, 1995.
  \item F. Buschmann et al., \textit{Pattern-Oriented Software Architecture}, Wiley, 1996.
  \item M. Fowler, \textit{Analysis Patterns}, Addison-Wesley, 1996.
  \item K. Beck, \textit{Smalltalk Best Practice Patterns}, Prentice Hall, 1997.
  \item S. Alpert et al., \textit{The Design Patterns Smalltalk Companion}, Addison-Wesley, 1998.
  \item W. Brown et al., \textit{AntiPatterns}, Wiley, 1998.
\end{itemize}

% ============================================================
\Needspace{30\baselineskip}
\gbox{In sintesi}{azure-gradient-6!80!black}{
\begin{itemize}
  \item Un \textbf{design pattern} è una soluzione comune a un problema ricorrente di progettazione, descritta da nome,
        problema, soluzione e conseguenze; nasce dall'architettura di Alexander e dal catalogo della \textbf{GoF} (1995).
  \item \textbf{Singleton}: costruttore privato, istanza in un campo statico, accesso con \code{getInstance()}; un solo
        oggetto, creato alla prima richiesta.
  \item \textbf{Builder}: la costruzione si fa a passi, su un oggetto builder (\code{with\ldots} restituisce \code{this}); il
        prodotto si crea, già validato, con \code{build()}. Il Director ripete la stessa ricetta con builder diversi.
  \item \textbf{Adapter}: una classe che implementa l'interfaccia attesa (Target) e traduce le chiamate verso la classe
        esistente (Adaptee), per ereditarietà o, più spesso, per composizione.
  \item \textbf{Composite}: foglie e composti hanno la stessa interfaccia; il composto contiene componenti, e le operazioni
        si propagano ricorsivamente sull'albero.
  \item \textbf{Decorator}: decoratori con la stessa interfaccia dell'oggetto decorato, concatenati a tempo di esecuzione;
        ogni chiamata attraversa la catena.
  \item \textbf{Observer}: il soggetto tiene la lista degli osservatori (\code{attach}, \code{detach}) e a ogni cambiamento li
        notifica; è il cuore dell'MVC e dei \textit{listener}.
  \item \textbf{State}: uno stato è un oggetto, una transizione è il cambio di un riferimento; \textbf{Strategy}: stessa
        struttura, ma è il client a scegliere l'algoritmo.
  \item \textbf{Dependency Injection}: le dipendenze si passano dall'esterno, di solito nel costruttore (inversione del
        controllo); così nei test si possono sostituire con oggetti fittizi.
\end{itemize}
}
